% GENERATED by tools/make_appendices.py from docs/extension/PHASE3_inference.md
% Do not edit by hand; edit the source and regenerate.
\emph{Builds on P1, 2a and 2b. Proof steps carry justification tags (2a §0.3, 2b §0.3, plus §0.3 below). Theorems marked ``(cited)'' apply a published result whose conditions are verified here but whose proof is not repeated. Checks: \path{verify_phase3.py} (17/17 PASS; \path{verify_phase3_output.txt}). Monte Carlo: \path{phase3_montecarlo.py}, 1000 replications per design, with Monte Carlo standard errors (\path{phase3_montecarlo_output.txt}; cross-check against the package in \path{phase3_montecarlo_crosscheck.py}).}

\begin{center}\rule{0.5\linewidth}{0.5pt}\end{center}

\hypertarget{pE-setting}{%
\subsection{0. Setting}\label{pE-setting}}

\hypertarget{pE-sampling}{%
\subsubsection{0.1 Sampling}\label{pE-sampling}}

\begin{itemize}
\tightlist
\item
  There are \(n\) units with observables \(W_i=(N_i,X_i,T_i,D_i)\).
\item
  \textbf{(P3-S1a)} Given the design and the district effects, the unit rates \(\beta_i\) are independent across units, with \(E[\beta_i\mid\text{design}]=m(V_i)\). Equivalently, the residuals \(e_i=\beta_i-m(V_i)\) are independent with mean zero. No dependence is allowed within a district beyond the district effect, and no clustered variance is used. §0.4 explains why. \emph{(Revised after review: the earlier P3-S1 allowed arbitrary within-district dependence for M0 while Theorem P3.5's residual allowance assumed independence.)}
\item
  \textbf{(P3-S2)} The design is conditioned on: inference is conditional on \(\{(N_i,X_i,D_i)\}\), so the vertices of the sample hulls \(\hat A_d\) are fixed.
\item
  \textbf{Regimes} (P1 A10): (R1) \(D^\star\) fixed, \(n_d\to\infty\); (R2) \(D^\star\to\infty\), \(n_d\) bounded; (R3) both grow. Which result holds where is in §0.5.
\end{itemize}

\hypertarget{pE-targets}{%
\subsubsection{0.2 Targets}\label{pE-targets}}

There are three related objects (P1 §1.7):

\begin{itemize}
\tightlist
\item
  \textbf{Superpopulation rate} \(B^*_{rc}=E[NX_r\beta_{rc}]/E[NX_r]\), the object of Phases 2a--2b.
\item
  \textbf{Finite-population rate} (JKST's target): \(B_{rc}=\sum_iN_iX_{ir}\beta_{irc}\big/\sum_iN_iX_{ir}\), the realised rate for this electorate.
\item
  \textbf{Finite-design conditional-mean rate:} \(\tilde B_{rc}=\sum_iw_{ir}\,m_{rc}(V_i)\), with \(w_{ir}=N_iX_{ir}/\sum_jN_jX_{jr}\) and \(m=E[\beta\mid V]\). This is the conditional expectation of \(B\) given the observed design (under P3-S1). It is a sample-weighted analogue of \(B^*\), not equal to it in general. Under suitable design-sampling conditions it estimates \(B^*\), but inference for \(B^*\) would also need the variation of the design, which is not pursued here.
\end{itemize}

All inference below is conditional on the design and covers \(B\) or \(\tilde B\).

They differ by the residual term
\[E_{rc}=B_{rc}-\tilde B_{rc}=\sum_iw_{ir}e_{irc},\qquad e=\beta-m.\]
\(\tilde B\) is the right target when the question is about behaviour rather than the particular electorate. Only \(\tilde B\) is affine in \(K\), with the \emph{sample} weights \(w\), so \(\mu\) is never estimated. \(B\) is \(\tilde B\) plus the residual term \(E\), which does not depend on \(K\) and is handled separately (§5). \emph{(Corrected after review: an earlier version said both are linear in \(K\).)}

\hypertarget{pE-additional-tags}{%
\subsubsection{0.3 Additional tags}\label{pE-additional-tags}}

\begin{longtable}[]{@{}
  >{\raggedright\arraybackslash\hspace{0pt}}p{(\columnwidth - 2\tabcolsep) * \real{0.5000}}
  >{\raggedright\arraybackslash\hspace{0pt}}p{(\columnwidth - 2\tabcolsep) * \real{0.5000}}@{}}
\toprule\noalign{}
\begin{minipage}[b]{\linewidth}\raggedright\hspace{0pt}Tag
\end{minipage} & \begin{minipage}[b]{\linewidth}\raggedright\hspace{0pt}Meaning
\end{minipage} \\
\midrule\noalign{}
\endhead
\bottomrule\noalign{}
\endlastfoot
{[}LPdual{]} & LP strong duality: a feasible LP with bounded value has an optimal dual solution with equal value; the optimum of a pointed polyhedron is attained at a vertex; a polyhedron has finitely many vertices \\
{[}Slater-b{]} & under strict feasibility, the optimal dual vectors are bounded: \(\lVert\lambda\rVert_1\le(\text{value}-\text{value at the strictly feasible point})/\varepsilon\) \\
{[}Hausdorff{]} & Hausdorff/Kuratowski convergence of compact convex sets implies convergence of support functions \\
{[}dDM{]} & delta method for Hadamard directionally differentiable maps (Shapiro 1991; Fang \& Santos 2019, Thm 2.1†) \\
{[}minLin{]} & the directional derivative of \(\beta\mapsto\min_{j\in J}g_j^\top\beta\) (\(J\) finite) at \(\beta_0\), in direction \(h\), is \(\min_{j\in J^*(\beta_0)}g_j^\top h\), with \(J^*\) the active set \\
{[}CLR{]} & Chernozhukov, Lee \& Rosen (2013, \emph{Econometrica} 81:667--737): parametric bounding functions (Condition P, Lemma 4, Thm 4, via Thms 1--2); finite index set as in their Example 1. \emph{Corrected after review: earlier cited as Thm 2; see 2a §0.4.} \\
{[}Pop{]} & Popoviciu: \(\operatorname{Var}(Y)\le(b-a)^2/4\) for \(Y\in[a,b]\); {[}BD{]} Bhatia--Davis: \(\operatorname{Var}(Y)\le(M-\mu)(\mu-m)\) \\
{[}Lyap{]} & Lindeberg--Feller CLT for independent summands \(Y_i\) with \(\lvert Y_i\rvert\le b_i\), applied relative to the \textbf{actual} standard deviation \(s_n\) of the sum: it holds when \(\max_ib_i/s_n\to0\). \emph{(Corrected after review, Batch B, Nov 2026: earlier stated as a CLT under \(\max_iw_i^2/\sum_iw_i^2\to0\) alone, which does not give a variance-standardised CLT when, e.g., only one summand has positive variance. P3.5(a) step 5 now combines it with Chebyshev; see there.)} \\
{[}union{]} & \(P(A\cup B)\le P(A)+P(B)\) \\
{[}OLS{]} & consistency and asymptotic normality of OLS for a fixed-dimensional parameter, with heteroskedasticity-robust (HC) variance, independent units, a full-rank design and bounded outcomes \\
{[}BR{]} & the multivariate CLT holds uniformly over all convex sets of \(\mathbb R^p\) (Bhattacharya \& Rao 1976, \emph{Normal Approximation and Asymptotic Expansions}, Thm 17.5, standard) \\
{[}Hoeff{]} & Hoeffding's inequality: for independent \(Y_i\in[a_i,b_i]\), \(P\big(\sum_i(Y_i-EY_i)\ge t\big)\le\exp\{-2t^2/\sum_i(b_i-a_i)^2\}\) (Hoeffding 1963, \emph{JASA} 58:13--30, Thm 2) \\
{[}CCK{]} & Gaussian approximation for maxima of high-dimensional sums (Chernozhukov, Chetverikov \& Kato 2017, \emph{Ann. Probab.} 45:2309--2352) \\
\end{longtable}

\hypertarget{pE-where-the-randomness-comes-from-added-after-review}{%
\subsubsection{0.4 Where the randomness comes from (added after review)}\label{pE-where-the-randomness-comes-from-added-after-review}}

Abadie, Athey, Imbens and Wooldridge (2020, \emph{Econometrica}; 2023, \emph{QJE}) distinguish two sources of uncertainty:

\begin{itemize}
\tightlist
\item
  \textbf{Sampling-based:} the units are a sample from a larger population.
\item
  \textbf{Design-based:} some assignment is random.
\end{itemize}

They show that clustered standard errors are justified by clustered sampling or clustered assignment, not by the mere fact that residuals are correlated. EI has neither. Every precinct and riding of the electorate is observed, and nothing is assigned. In a purely design-based view, nothing is random at all: \(B\) is a fixed unknown, and the only valid bound is Duncan--Davis.

The contextual-effects assumption is a statement about a conditional expectation, \(E[\beta\mid V]=M_H(X)\). That only has content if the precinct rates are treated as random given the precinct's characteristics. The uncertainty here is therefore \textbf{model-based, conditional on the design}: the \(\beta_i\) are random given \(V_i\), and the design is held fixed (P3-S2). Abadie, Imbens and Zheng (2014) give the parallel fixed-regressor framework for regression.

Within this framework:

\begin{itemize}
\tightlist
\item
  \textbf{District shocks are fixed features of this electorate.} The targets \(B\) and \(\tilde B\) refer to this electorate, so a shock common to a riding (candidate, campaign, local swing) is not re-drawn.

  \begin{itemize}
  \tightlist
  \item
    Under M1/M2 it is a district effect, part of the mean \(m(V_i)\).
  \item
    Under M0 a district shock is not a variance problem but a misspecification. Conditional on the design, M0's mean model holds for the observed units only if the district shifts are zero (Prop. P3.6(c)). Clustered standard errors would not repair it.
  \end{itemize}
\item
  \textbf{What remains assumed} (P3-S1a): no dependence within a district beyond its district effect. Neighbourhood spillovers between adjacent precincts would violate it. This is a stated limitation. A partial check is residual correlation between neighbouring precincts after fitting M1.
\item
  \textbf{Independence is a substantive assumption} \emph{(added after review, round 2, Nov 2026)}. Two arguments must be kept apart.

  \begin{enumerate}
  \def\labelenumi{\arabic{enumi}.}
  \tightlist
  \item
    Conditioning on observed district effects and \textbf{assuming} that the remaining residuals are independent (P3-S1a) is a coherent model.
  \item
    ``Every unit is observed, so dependence need not be considered'' does \textbf{not} follow. Observing every unit removes sampling uncertainty about which units exist; it says nothing about dependence among their residuals.
  \end{enumerate}

  Within-district shocks can make \(\operatorname{Cov}(e_i,e_j\mid\text{design})\ne0\) with no district mean shift. That violates P3-S1a without violating the mean model, so the pooling test P3.6(c), which tests mean shifts, is \textbf{not} a test of within-district dependence. P3-S1a must therefore be argued for each application; residual correlation within districts and between neighbours is a diagnostic, not a test that can confirm it. We do not reintroduce clustered standard errors, for the reasons above. Under dependence, the variance and residual-allowance statements of P3.5 and P3.8 are not covered.
\item
  \textbf{A superpopulation target that averages over re-drawn district shocks} (``behaviour in a typical riding'') would need cluster methods and many districts. It is not pursued.
\end{itemize}

\hypertarget{pE-which-result-holds-in-which-regime}{%
\subsubsection{0.5 Which result holds in which regime}\label{pE-which-result-holds-in-which-regime}}

\begin{longtable}[]{@{}
  >{\raggedright\arraybackslash\hspace{0pt}}p{(\columnwidth - 6\tabcolsep) * \real{0.2500}}
  >{\raggedright\arraybackslash\hspace{0pt}}p{(\columnwidth - 6\tabcolsep) * \real{0.2500}}
  >{\raggedright\arraybackslash\hspace{0pt}}p{(\columnwidth - 6\tabcolsep) * \real{0.2500}}
  >{\raggedright\arraybackslash\hspace{0pt}}p{(\columnwidth - 6\tabcolsep) * \real{0.2500}}@{}}
\toprule\noalign{}
\begin{minipage}[b]{\linewidth}\raggedright\hspace{0pt}Model
\end{minipage} & \begin{minipage}[b]{\linewidth}\raggedright\hspace{0pt}Regime
\end{minipage} & \begin{minipage}[b]{\linewidth}\raggedright\hspace{0pt}Formal result
\end{minipage} & \begin{minipage}[b]{\linewidth}\raggedright\hspace{0pt}Status
\end{minipage} \\
\midrule\noalign{}
\endhead
\bottomrule\noalign{}
\endlastfoot
any rung & any (R1, R2, R3), any \(n\) & Theorem P3.8: finite-sample CI; finite-sample emptiness test & Theorem; currently uninformative in width (§5.1) \\
M0 & any (\(\dim\beta\) fixed) & P3.2, P3.4, P3.5(a), (b1), (b2) & Theorem (asymptotic, \(n\to\infty\)) \\
M1, M2 & R1 (\(D^\star\) fixed) & P3.2, P3.4, P3.5(a), (b1), (b2) & Theorem (asymptotic) \\
M1, M2 & R3 & the same, with a high-dimensional Gaussian approximation {[}CCK{]} under a rate condition on \(D^\star\) versus \(n_d\) & Conjecture \\
M1 & R2 & coverage of the CLR interval & Simulation finding (§8.2: 0.995--1 with up to 200 districts of 10 polls) \\
any & any & the CLR interval as implemented (candidates by simulation) & Conjecture P3.5(b3) \\
\end{longtable}

Under M0 the number of constraints grows with the number of hull vertices, but the regression has fixed dimension \(p\). Every constraint term is a linear functional of one \(p\)-dimensional asymptotically normal vector. The relevant events, \(\{\max_j(\cdot)\le k\}\), are convex sets in \(\mathbb R^p\), and the multivariate CLT holds uniformly over convex sets {[}BR{]}. So the growth in constraints does not matter. Under M1/M2 in R2 or R3, \(\dim\beta\) grows with \(D^\star\), and that argument no longer applies.

\begin{center}\rule{0.5\linewidth}{0.5pt}\end{center}

\hypertarget{pE-estimators}{%
\subsection{1. Estimators}\label{pE-estimators}}

\begin{enumerate}
\def\labelenumi{\arabic{enumi}.}
\tightlist
\item
  \textbf{Regression.} Fit each column \(T_c\) by OLS on the homogeneous quadratic monomials \(X_rX_s\) (\(r\le s\)) for M0. For M1, also include the district shifts \(1\{D=d\}X_r\), \(d\ge2\); this imposes the overidentifying restrictions of 2b Lemma D1(b). For M2, fit one regression per district. The coefficients are \(\hat\beta\), and the estimated columns are coherent automatically (2a Remark 1.2). Use an HC3 covariance \(\hat\Sigma\) (P3-S1a; no clustering, §0.4). Only \(C-1\) columns are needed; column \(C\) follows by adding up (Remark P3.4.1).
\item
  \textbf{Constraint terms.} For each district \(d\), vertex \(v\in\operatorname{vert}(\hat A_d)\), row \(r\) and column \(c\), the constraint is \(s_j(\beta)+a_j^\top k\ge0\), with \(s_j\) linear in \(\beta\) and \(a_j\) fixed (2a Prop. 7, 2b Thm 2(b)).
\item
  \textbf{Target.} \(\tilde B(k)=c(\beta)+\gamma^\top k\), with \(c\) linear in \(\beta\) and \(\gamma\) computed from the sample weights (2a Lemma 2, 2b Lemma D2).
\item
  \textbf{Plug-in bound.} \([\hat\varphi_L,\hat\varphi_U]=c(\hat\beta)+[\min,\max]\{\gamma^\top k:s(\hat\beta)+Ak\ge0\}\).
\end{enumerate}

\begin{center}\rule{0.5\linewidth}{0.5pt}\end{center}

\hypertarget{pE-structure-of-the-bound}{%
\subsection{2. Structure of the bound}\label{pE-structure-of-the-bound}}

\phantomsection\label{pE-res-lemma-P3.1}\textbf{Lemma P3.1 (the endpoint is a minimum of finitely many linear functions).} Suppose the LP is feasible with bounded value at \(\beta\). Then
\[\varphi_U(\beta)=c(\beta)+\min_{\lambda\in\Lambda}\lambda^\top s(\beta),\qquad\Lambda=\operatorname{vert}\{\lambda\ge0:A^\top\lambda=-\gamma\},\]
a finite set. \(\varphi_U\) is therefore concave and piecewise linear in \(\beta\), and likewise \(\varphi_L(\beta)=c(\beta)-\min_{\lambda\in\Lambda'}\lambda^\top s(\beta)\), with \(\Lambda'=\operatorname{vert}\{\lambda\ge0:A^\top\lambda=\gamma\}\).

\emph{Proof.}

\begin{enumerate}
\def\labelenumi{\arabic{enumi}.}
\tightlist
\item
  \(\max\{\gamma^\top k:s+Ak\ge0\}\) has Lagrangian \(\gamma^\top k+\lambda^\top(s+Ak)\). Stationarity in the free \(k\) requires \(\gamma+A^\top\lambda=0\), and the dual is \(\min\{\lambda^\top s:\lambda\ge0,\ A^\top\lambda=-\gamma\}\). {[}LPdual{]}
\item
  Primal feasibility and boundedness give equal optimal values. {[}LPdual; 2a:L3(b){]}
\item
  The dual feasible set lies in \(\{\lambda\ge0\}\), so it contains no line and is pointed. A linear objective bounded below on it attains its minimum at a vertex, and there are finitely many vertices. {[}LPdual{]}
\item
  The dual feasible set does not depend on \(\beta\) (\(A\) and \(\gamma\) are fixed given the design). So \(\varphi_U(\beta)-c(\beta)=\min_{\lambda\in\Lambda}\lambda^\top s(\beta)\) for every \(\beta\) at which the primal is feasible and bounded. {[}3; Def{]}
\item
  A minimum of finitely many linear functions is concave and piecewise linear. {[}alg{]} The lower endpoint follows by applying 1--5 to \(-\gamma\). ∎
\end{enumerate}

Checked: concavity and piecewise linearity along random directions (\path{verify_phase3.py}).

\emph{Scope.} The lemma is about variation in the regression coefficients \(\beta\) for a \textbf{fixed} constraint matrix \(A\) and objective \(\gamma\). Conditioning on the design fixes both, because the target uses the sample weights. It is the first case of 2a Prop. 7's regularity paragraph. It says nothing about joint variation in the identified coefficients, the population moments and the constraint set, where the value function is in general neither concave nor piecewise linear.

\textbf{Reading.} JKST's Proposition 4 is this lemma with \(|\Lambda|=4\) (their \(g_{lj},g_{uj}\)). Its ``tie-breaking'' assumption says the minimising \(\lambda\) is unique. In R×C, and especially with many districts, \(\Lambda\) is large and near-ties are the norm. That is why the inference below treats the endpoint as an intersection bound.

\begin{center}\rule{0.5\linewidth}{0.5pt}\end{center}

\hypertarget{pE-consistency-and-rate}{%
\subsection{3. Consistency and rate}\label{pE-consistency-and-rate}}

\phantomsection\label{pE-res-proposition-P3.2}\textbf{Proposition P3.2.} Assume (P3-S1--S2), {[}A1{]}--{[}A5{]} for the rung used, \(\sqrt n(\hat\beta-\beta_0)=O_p(1)\), and \textbf{(Slater)}: there are \(\tilde k\) and \(\varepsilon>0\) with \(s_j(\beta_0)+a_j^\top\tilde k\ge\varepsilon\) for all constraints defined by the population hulls. Then:

\begin{enumerate}
\def\labelenumi{(\alph{enumi})}
\tightlist
\item
  \(\hat\varphi_U-\varphi_U(\beta_0;\hat A)=O_p(n^{-1/2})\), and likewise for \(\varphi_L\);
\item
  \(\varphi_U(\beta_0;\hat A)\to\varphi_U(\beta_0;A)\) when \(\hat A\to A\) in Hausdorff distance, either along a nested sequence (\(\hat A_n\subseteq\hat A_{n+1}\), as when observations are added), or for any Hausdorff-convergent sequence whose feasible sets \(\Theta(\hat A_n)\) eventually lie in a common bounded set;
\item
  since \(\hat A_d\subseteq A_d\), \(\varphi_U(\beta_0;\hat A)\ge\varphi_U(\beta_0;A)\) and \(\varphi_L(\beta_0;\hat A)\le\varphi_L(\beta_0;A)\): \textbf{estimating the hull only ever widens the bound.}
\end{enumerate}

\emph{Proof.}

\begin{enumerate}
\def\labelenumi{\arabic{enumi}.}
\tightlist
\item
  (c): a vertex of \(\hat A_d\) is an observed composition, hence in \(A_d\). So \(\hat A_d\subseteq A_d\), the constraints over \(\hat A_d\) are a subset of those over \(A_d\), and the feasible set is weakly larger. {[}Def; 2a:L3(a){]}
\item
  (a): by Lemma P3.1, \(\varphi_U(\beta)-\varphi_U(\beta')=c(\beta-\beta')+\big[\min_\lambda\lambda^\top s(\beta)-\min_\lambda\lambda^\top s(\beta')\big]\). {[}P3.1{]}
\item
  The bracket is bounded by \(\max\lVert S^\top\lambda\rVert\,\lVert\beta-\beta'\rVert\), the max taken over optimal duals at \(\beta\) or \(\beta'\). {[}P3.1; \(s\) linear, \(s(\beta)=S\beta\){]}
\item
  For \(\beta\) in a neighbourhood of \(\beta_0\), \(\tilde k\) remains \(\varepsilon/2\)-strictly feasible, so the optimal duals are uniformly bounded. {[}Slater; continuity of \(s\); Slater-b{]}
\item
  Hence \(\lvert\hat\varphi_U-\varphi_U(\beta_0;\hat A)\rvert\le L\lVert\hat\beta-\beta_0\rVert=O_p(n^{-1/2})\), with \(L\) uniform in the hull because the constraint set over \(\hat A\) is a subset of the one over \(A\). {[}2--4; 1{]}
\item
  (b): the sets \(\Theta(\hat A)\) are compact and convex, decrease as \(\hat A\) grows, and contain \(\Theta(A)\). {[}1; 2a:L3(b){]}
\item
  If \(k\in\Theta(\hat A_p)\) for all \(p\) and \(v\in A\), take \(v_p\in\hat A_p\) with \(v_p\to v\). The constraints are continuous in \(v\), so \(k\) satisfies the constraint at \(v\). Hence \(\bigcap_p\Theta(\hat A_p)=\Theta(A)\). {[}Def; continuity{]}
\item
  A decreasing sequence of compact convex sets converges in Hausdorff distance to its intersection, so the support functions, and hence the LP values, converge. {[}Hausdorff{]} ∎(nested)
\item
  Non-nested case: \(\Theta(A)\subseteq\Theta(\hat A_n)\) for every \(n\) by 1. If \(k_n\in\Theta(\hat A_n)\) and \(k_{n_j}\to k\), then for each \(v\in A\) pick \(v_j\in\hat A_{n_j}\) with \(v_j\to v\). Continuity of the constraints in \((k,v)\) gives the constraint at \(v\), so \(k\in\Theta(A)\). The upper Kuratowski limit is therefore \(\Theta(A)\), which is also contained in every term. In a common bounded set this is Hausdorff convergence. {[}1; Def; continuity; Hausdorff{]} ∎(b)
\end{enumerate}

\textbf{Remarks.}

\begin{itemize}
\tightlist
\item
  \textbf{Error decomposition.} \(\hat\varphi_U-\varphi_U(\beta_0;A)=\underbrace{[\hat\varphi_U-\varphi_U(\beta_0;\hat A)]}_{O_p(n^{-1/2})\text{ by (a)}}+\underbrace{[\varphi_U(\beta_0;\hat A)-\varphi_U(\beta_0;A)]}_{\ge0,\ \to0\text{ by (b), rate not shown}}\). The second term is not shown to be \(O_p(n^{-1/2})\), and can be larger. So Theorem P3.4's limit concerns the \textbf{sample-hull} endpoint \(\varphi_U(\beta_0;\hat A)\), and not the population-hull endpoint. Theorem P3.4 itself is for a \textbf{fixed} constraint system; as the sample hull changes with \(n\), the limit is given by Proposition P3.4′ under its localisation condition. \emph{(Corrected after review, round 2, Nov 2026.)}
\item
  \textbf{Slater is a strict-interiority condition.} It requires a \(k\) at which every implied rate at every hull vertex is at least \(\varepsilon>0\). It fails when every feasible \(k\) makes some implied rate exactly 0 at some vertex, e.g.~a true structural zero, which is not exotic in EI. Design E of §8 is an example: the identified set is a single point. (Design C has a rate close to 0 at one vertex: Slater holds, but only just.) The Lipschitz bound in (a) is proved only under Slater; boundary cases are open (§10). The coverage results of §5 do \textbf{not} use Slater: P3.5(a) step 2 uses only that the true \(k_0\) is feasible.
\item
  The hull effect in (c) is a one-sided (conservative) bias. Its size depends on how much mass \(P_X\) puts near the boundary of \(A_d\); it can vanish slowly. Report the bound's sensitivity to trimming the hull (P1 A4).
\item
  The \(O_p(n^{-1/2})\) rate is for the estimation part. The target \(B\) adds \(E=O_p(n^{-1/2})\) as well (§5). Bounds do not shrink to a point: their width converges to the population width, which is \(O(1)\).
\end{itemize}

\begin{center}\rule{0.5\linewidth}{0.5pt}\end{center}

\hypertarget{pE-asymptotic-distribution}{%
\subsection{4. Asymptotic distribution}\label{pE-asymptotic-distribution}}

\phantomsection\label{pE-res-theorem-P3.4}\textbf{Theorem P3.4 (cited structure: {[}dDM{]}; fixed constraint system).} Under Proposition P3.2's assumptions and \(\sqrt n(\hat\beta-\beta_0)\Rightarrow\mathcal Z\sim N(0,\Sigma)\), suppose the \textbf{constraint system is fixed} along the sequence: the vertex set entering the LP does not depend on \(n\), so \(A\), \(\gamma\), \(\Lambda\) and the gradients \(g_\lambda\) are fixed. This holds for a prespecified polytope \(A\) with fixed vertices, or along a design sequence whose sample hull is eventually constant, \(\hat A_n=A\) for \(n\ge n_0\) (for instance, when \(A\) is a polytope whose vertices are atoms of the design law). Then
\[\sqrt n\big(\hat\varphi_U-\varphi_U(\beta_0)\big)\Rightarrow\min_{\lambda\in\Lambda^*}\big(c+S^\top\lambda\big)^\top\mathcal Z,\]
where \(\Lambda^*\subseteq\Lambda\) is the set of optimal dual vertices at \(\beta_0\).

\begin{itemize}
\item
  \begin{enumerate}
  \def\labelenumi{(\roman{enumi})}
  \tightlist
  \item
    If \(\Lambda^*\) is a singleton (JKST's tie-breaking), the limit is \(N(0,g^{*\top}\Sigma g^*)\) with \(g^*=c+S^\top\lambda^*\).
  \end{enumerate}
\item
  \begin{enumerate}
  \def\labelenumi{(\roman{enumi})}
  \setcounter{enumi}{1}
  \tightlist
  \item
    If two active gradients differ in a direction in which the estimator varies, i.e.~\((g_\lambda-g_{\lambda'})^\top\Sigma(g_\lambda-g_{\lambda'})>0\) for some \(\lambda,\lambda'\in\Lambda^*\), the limit is a minimum of correlated normals that are not a.s. equal, and it has strictly negative mean. The plug-in upper endpoint is then biased \emph{downward}, and the lower endpoint upward, both by order \(n^{-1/2}\). That is, the plug-in bound is too narrow. If instead all active gradients agree on the support of \(\mathcal Z\), all the \(g_\lambda^\top\mathcal Z\) are a.s. equal and the limit is normal, as in (i), even though \(\Lambda^*\) is not a singleton. \emph{(Corrected after review: an earlier version claimed a negative mean whenever \(\Lambda^*\) is not a singleton.)}
  \end{enumerate}
\item
  \begin{enumerate}
  \def\labelenumi{(\roman{enumi})}
  \setcounter{enumi}{2}
  \tightlist
  \item
    In case (ii) the standard bootstrap is inconsistent (Fang \& Santos 2019, Cor. 3.1†).
  \end{enumerate}
\end{itemize}

\emph{(Corrected after review, round 2, Nov 2026. The theorem was stated ``with the design and the sample hull \(\hat A\) fixed''. The proof is correct for a fixed set of affine maps, but as \(n\) grows the sample hull, the dual vertices and the active set can change, and gaps between nearly binding candidates can shrink at rate \(n^{-1/2}\). The referee's example: \(f_n(\theta)=\min\{\theta,-\theta+n^{-1/2}\}\), \(\theta_0=0\), \(\sqrt n(\hat\theta-\theta_0)\Rightarrow Z\). The first piece is uniquely active for every \(n\), yet \(\sqrt n(f_n(\hat\theta)-f_n(\theta_0))\Rightarrow\min\{Z,1-Z\}\), which is not normal. The theorem is now stated for a fixed constraint system; changing systems are Proposition P3.4′. Neither affects the coverage results of §5, which do not use P3.4.)}

\emph{Proof.}

\begin{enumerate}
\def\labelenumi{\arabic{enumi}.}
\tightlist
\item
  \(\varphi_U(\beta)=\min_{\lambda\in\Lambda}g_\lambda^\top\beta\) with \(g_\lambda=c+S^\top\lambda\): a finite minimum of linear maps. {[}P3.1{]}
\item
  Such a map is Lipschitz, and its directional derivative is \(h\mapsto\min_{\lambda\in\Lambda^*}g_\lambda^\top h\). A Lipschitz map with directional derivatives in finite dimensions is Hadamard directionally differentiable. {[}minLin{]}
\item
  Apply the directional delta method. {[}dDM; 1--2{]}
\item
  (i): the minimum over a singleton is linear, giving a normal limit. {[}3{]}
\item
  (ii): let \(Y=g_\lambda^\top\mathcal Z\) and \(Y'=g_{\lambda'}^\top\mathcal Z\) with \(\operatorname{Var}(Y-Y')=(g_\lambda-g_{\lambda'})^\top\Sigma(g_\lambda-g_{\lambda'})>0\). Then \(E\min_{\Lambda^*}g^\top\mathcal Z\le E\min(Y,Y')=\tfrac12E(Y+Y')-\tfrac12E\lvert Y-Y'\rvert=-\tfrac12E\lvert Y-Y'\rvert<0\). If \(\operatorname{Var}(g_\lambda^\top\mathcal Z-g_{\lambda'}^\top\mathcal Z)=0\) for all active pairs, the minimum equals each \(g_\lambda^\top\mathcal Z\) a.s., which is normal. {[}3; \(\min(a,b)=\frac{a+b}2-\frac{\lvert a-b\rvert}2\); \(Y-Y'\) is a nondegenerate normal{]}
\item
  (iii): in case (ii) the derivative is nonlinear on the support of \(\mathcal Z\), so the condition for bootstrap consistency fails. {[}Fang \& Santos 2019†{]} ∎
\end{enumerate}

\phantomsection\label{pE-res-remark-P3.4.1}\textbf{Remark P3.4.1 (parametrise \(C-1\) columns).} The shares satisfy \(\sum_cT_c=1\), so the coefficient vectors of the \(C\) column regressions sum to a known constant, and the joint covariance of all \(C\) columns is singular. State \(\hat\beta\), \(\Sigma\) and the gradients \(g_\lambda\) for the first \(C-1\) columns and recover column \(C\) by adding up. Then \(\Sigma\) can be nonsingular. When it is positive definite, the condition in (ii) reduces to \(g_\lambda\neq g_{\lambda'}\) in the reduced parametrisation.

\phantomsection\label{pE-res-proposition-P3.4p}\textbf{Proposition P3.4′ (triangular arrays; added after review, round 2).} Let the design, and hence the constraint system, change with \(n\). Write \(\Lambda_n\) for the dual vertex set of the design-\(n\) LP, with \(\lvert\Lambda_n\rvert\le\bar L\) for all \(n\), gradients \(g_{\lambda,n}=c_n+S_n^\top\lambda\) (bounded), and \textbf{localised slacks}
\[c_{\lambda,n}:=\sqrt n\big(\theta_{\lambda,n}(\beta_0)-\varphi_{U,n}(\beta_0)\big)\ \ge0,\qquad\theta_{\lambda,n}(\beta)=g_{\lambda,n}^\top\beta .\]
\textbf{Localisation condition.} After relabelling the candidates, \(g_{\lambda,n}\to g_\lambda\) and \(c_{\lambda,n}\to c_\lambda\in[0,\infty]\) for each \(\lambda\). Then
\[\sqrt n\big(\hat\varphi_{U,n}-\varphi_{U,n}(\beta_0)\big)\Rightarrow\min_{\lambda:\ c_\lambda<\infty}\big(g_\lambda^\top\mathcal Z+c_\lambda\big).\]
(Without the condition, the same holds along every subsequence on which it holds; the limit can then differ between subsequences.)

\begin{itemize}
\tightlist
\item
  \textbf{Separation condition.} The limit of Theorem P3.4 (a minimum of \(g_\lambda^\top\mathcal Z\) over the active set) obtains when every \(c_\lambda\in\{0,\infty\}\): each candidate is either exactly active at \(\beta_0\) or inactive by more than \(n^{-1/2}\), \(\sqrt n\min_{\lambda\notin\Lambda^*_n}\{\theta_{\lambda,n}(\beta_0)-\varphi_{U,n}(\beta_0)\}\to\infty\), with a stable active set and convergent gradients.
\item
  A candidate inactive by a gap of exact order \(n^{-1/2}\) (\(0<c_\lambda<\infty\)) stays in the limit. The referee's \(f_n\) has \(c=1\) for the second piece and gives \(\min\{Z,1-Z\}\).
\item
  This is a limit for the sample-hull endpoint along the design sequence. It says nothing about the population-hull error (remark after P3.2).
\end{itemize}

\emph{Proof.}

\begin{enumerate}
\def\labelenumi{\arabic{enumi}.}
\tightlist
\item
  \(\varphi_{U,n}(\beta)=\min_{\lambda\in\Lambda_n}\theta_{\lambda,n}(\beta)\) for every \(\beta\) at which the LP is feasible and bounded. {[}P3.1, for the design-\(n\) system{]}
\item
  Hence, exactly, \(\sqrt n(\hat\varphi_{U,n}-\varphi_{U,n}(\beta_0))=\min_\lambda\big(g_{\lambda,n}^\top\sqrt n(\hat\beta-\beta_0)+c_{\lambda,n}\big)\). {[}1; Def \(c_{\lambda,n}\); linearity{]}
\item
  \(\sqrt n(\hat\beta-\beta_0)=O_p(1)\) and the gradients are bounded, so a candidate with \(c_{\lambda,n}\to\infty\) exceeds the active candidates' terms with probability \(\to1\) and drops out of the minimum. {[}OLS; 2{]}
\item
  The remaining finitely many terms converge jointly in distribution to \((g_\lambda^\top\mathcal Z+c_\lambda)_\lambda\), and the minimum is continuous. {[}2--3; continuous mapping{]} ∎
\end{enumerate}

Checked (\path{verify_phase3.py}, ``P3.4′ {[}F{]}''): the referee's \(f_n\) at \(n=10^4\); the limit is \(\min\{Z,1-Z\}\) (mean \(\approx-0.40\), Kolmogorov distance 0.31 from \(N(0,1)\)), not the normal limit the old reading gives.

Checked: with two tied constraints and unit noise, the mean plug-in error is \(-0.564\approx-1/\sqrt\pi\); with a unique active constraint it is 0.007. With two active gradients that differ only in a null direction of \(\Sigma\), the mean is 0 (normal limit), as the corrected (ii) states.

\textbf{Boundary cases.}

\begin{itemize}
\tightlist
\item
  \textbf{Truth at a vertex.} When the true \(K\) lies on the boundary of the identified set, the side of the interval that touches it carries the full error probability, so coverage is close to nominal. When the true \(K\) is interior, coverage tends to 1 (JKST App. A Remark). This is the Imbens--Manski distinction between covering the parameter and covering the set.
\item
  \textbf{Width near zero.} When the bound's width is near zero (point identification: L4 in 2×2, L5, the DiD with unchanged contextual effects), one-sided intervals per endpoint under-cover unless the critical value adapts (Stoye 2009). The two-sided Bonferroni split used below is uniformly valid, at some cost in width.
\end{itemize}

\begin{center}\rule{0.5\linewidth}{0.5pt}\end{center}

\hypertarget{pE-confidence-intervals}{%
\subsection{5. Confidence intervals}\label{pE-confidence-intervals}}

\phantomsection\label{pE-res-theorem-P3.5}\textbf{Theorem P3.5 (validity).} Fix \(\alpha_U,\alpha_L,\alpha_E\in(0,1)\). Write \(\theta_\lambda(\beta)=g_\lambda^\top\beta\) for the candidate upper endpoints (\(\lambda\in\Lambda\)), with standard errors \(\hat\sigma_\lambda\) from \(\hat\Sigma\).

\textbf{(a) Conservative version.} Let \(\hat k\) be the \((1-\alpha_U-\alpha_L)\)-quantile of \(\max_j\) over all constraint terms and the constant of the standardised Gaussian errors (one-sided for the constraints, two-sided for the constant). Put
\[\hat\Theta^+=\{k:s(\hat\beta)+\hat k\hat\sigma_s+Ak\ge0\},\]
\[\resizebox{\linewidth}{!}{$\displaystyle \mathrm{CI}=\Big[c(\hat\beta)-\hat k\hat\sigma_c+\min_{\hat\Theta^+}\gamma^\top k-z_{\alpha_E}S_1,\ \ c(\hat\beta)+\hat k\hat\sigma_c+\max_{\hat\Theta^+}\gamma^\top k+z_{\alpha_E}S_1\Big],$}\]
with \(S_1=\tfrac12\big(\sum_iw_{ir}^2\big)^{1/2}\) and \(z_{\alpha_E}=\Phi^{-1}(1-\alpha_E/2)\). Then
\[\liminf_nP(B_{rc}\in\mathrm{CI})\ge1-\alpha_U-\alpha_L-\alpha_E.\]

\textbf{(b1) CLR with the full candidate set (cited: {[}CLR{]}, Thm 4).} Let \(\Lambda\) be the full set of dual vertices of Lemma P3.1, which is finite for a fixed design. Assume CLR's Condition P: \(\Sigma\) is positive definite (use \(C-1\) columns, Remark P3.4.1, and a full-rank design), and \(g_\lambda\ne0\) for every \(\lambda\in\Lambda\). For the upper end:

\begin{enumerate}
\def\labelenumi{\arabic{enumi}.}
\tightlist
\item
  Select \(\hat V=\{\lambda:\hat\theta_\lambda\le\min_\mu(\hat\theta_\mu+k_p\hat\sigma_\mu)+2k_p\hat\sigma_\lambda\}\), where \(k_p\) is the \(\gamma_n=1-0.1/\log n\) quantile of \(\max_\lambda Z_\lambda\), and \(Z\) is the standardised Gaussian vector with correlation from \(\hat\Sigma\).
\item
  Set \(\hat U=\min_{\lambda\in\hat V}(\hat\theta_\lambda+\hat k\hat\sigma_\lambda)\), with \(\hat k\) the \((1-\alpha_U)\)-quantile of \(\max_{\lambda\in\hat V}Z_\lambda\).
\item
  Do the same for the lower end, and add \(z_{\alpha_E}S_1\) on each side (each side's E-allowance at level \(\alpha_E/2\)).
\end{enumerate}

The same coverage bound holds. With \(\alpha_E=0\) and no \(S_1\) term, either version covers \(\tilde B_{rc}\) with probability at least \(1-\alpha_U-\alpha_L\).

\textbf{(b2) Any candidate set, with a simultaneous critical value (proved here).} Let \(\hat\Lambda\) be any data-dependent set of dual-feasible points (vertices or not; for example collected by simulation), and \(\hat U=\min_{\lambda\in\hat\Lambda}(\hat\theta_\lambda+\hat k\hat\sigma_\lambda)\), with \(\hat U=+\infty\) if \(\hat\Lambda=\emptyset\). Then \(\liminf_nP(\varphi_U(\beta_0)\le\hat U)\ge1-\alpha_U\) in either of two cases:

\begin{itemize}
\item
  \begin{enumerate}
  \def\labelenumi{(\roman{enumi})}
  \tightlist
  \item
    \(\hat k\) is the \((1-\alpha_U)\)-quantile of \(\max_{\lambda\in\bar\Lambda}Z_\lambda\) for a fixed finite set \(\bar\Lambda\supseteq\hat\Lambda\) a.s. (e.g.~all of \(\Lambda\)), with \(\sigma_\lambda>0\) on \(\bar\Lambda\);
  \end{enumerate}
\item
  \begin{enumerate}
  \def\labelenumi{(\roman{enumi})}
  \setcounter{enumi}{1}
  \tightlist
  \item
    projection: \(\hat k=\sqrt{\chi^2_{q,1-\alpha}}\) and \(\hat\sigma_\lambda=(g_\lambda^\top\hat\Sigma g_\lambda/n)^{1/2}\), in either of two settings: (ii-a) \(\Sigma\) is positive definite and \(q=\dim\beta\) (use \(C-1\) columns and a full-rank design, Remark P3.4.1); or (ii-b) a \textbf{reduced full-rank parameterisation}: there are a fixed matrix \(L\) (\(\dim\beta\times q\)) and a vector \(b\) with \(\beta=L\eta+b\) and \(\hat\beta=L\hat\eta+b\) \textbf{exactly} (so every error in a null direction of \(\Sigma\) is identically zero, e.g.~coordinates fixed by known identities), \(\sqrt n(\hat\eta-\eta_0)\Rightarrow N(0,\Sigma_\eta)\) with \(\Sigma_\eta\) positive definite, and \(\hat\Sigma=L\hat\Sigma_\eta L^\top\) with \(\hat\Sigma_\eta\) consistent. This one event covers both ends, any \(\hat\Lambda\), at total level \(\alpha\).
  \end{enumerate}

  \emph{(Corrected after review, Batch B, Nov 2026. The earlier statement took \(q=\operatorname{rank}\Sigma\) for any \(\Sigma\), with a pseudoinverse in the proof. That bounds \(g^\top(\hat\beta-\beta_0)\) only for \(g\) in the range of \(\Sigma\). A singular limit does not control errors of smaller order in its null directions: with \(\hat\beta-\beta_0=(Z/\sqrt n,\,-1/n)\), \(Z\sim N(0,1)\), the limit covariance is \(\operatorname{diag}(1,0)\), the second coordinate gets standard error 0, and the upper bound \(\hat\beta_2+0\) misses \(\beta_{0,2}\) in every sample. Counterexample due to the referee; \path{verify_phase3.py}, ``P3.5(b2)(ii) null direction''.)}
\end{itemize}

\emph{Proof.}

\begin{enumerate}
\def\labelenumi{\arabic{enumi}.}
\tightlist
\item
  Under the model the primal LP at \(\beta_0\) is feasible (\(k_0\) is feasible). For any dual-feasible \(\lambda\) and any primal-feasible \(k\): \(\gamma^\top k=-\lambda^\top Ak\le\lambda^\top s(\beta_0)\). So \(\theta_\lambda(\beta_0)\ge\varphi_U(\beta_0)\) for \textbf{every} dual-feasible \(\lambda\), not only optimal ones. {[}weak duality; \(\lambda\ge0\), \(s+Ak\ge0\){]}
\item
  On \(\mathcal E=\{\theta_\lambda(\beta_0)\le\hat\theta_\lambda+\hat k\hat\sigma_\lambda\ \forall\lambda\in\bar\Lambda\}\): \(\hat U\ge\min_{\lambda\in\hat\Lambda}\theta_\lambda(\beta_0)\ge\varphi_U(\beta_0)\). {[}1; \(\hat\Lambda\subseteq\bar\Lambda\){]}
\item
  (i): \((\theta_\lambda(\beta_0)-\hat\theta_\lambda)/\hat\sigma_\lambda\), \(\lambda\in\bar\Lambda\), is a fixed finite collection of standardised linear functionals of \(\hat\beta\), jointly asymptotically \(N(0,\text{corr})\). So \(P(\mathcal E)\to1-\alpha_U\). {[}OLS; \(\hat\Sigma\) consistent; Def \(\hat k\){]}
\item
  (ii): case (ii-b) contains (ii-a) (take \(L=I\), \(b=0\), \(\eta=\beta\)). Write \(\hat e=\hat\eta-\eta_0\). For every \(g\), \(g^\top(\hat\beta-\beta_0)=(L^\top g)^\top\hat e\) exactly, so by Cauchy--Schwarz in the metric of \(\hat\Sigma_\eta\) (positive definite with probability \(\to1\)),
  \[\resizebox{\linewidth}{!}{$\displaystyle \lvert g^\top(\hat\beta-\beta_0)\rvert\le\big(\hat e^\top\hat\Sigma_\eta^{-1}\hat e\big)^{1/2}\big(g^\top L\hat\Sigma_\eta L^\top g\big)^{1/2}=\big(n\,\hat e^\top\hat\Sigma_\eta^{-1}\hat e\big)^{1/2}\hat\sigma_g .$}\]
  This holds for every \(g\) at once, including \(g\) with \(L^\top g=0\), for which both sides are 0. By continuous mapping, \(n\,\hat e^\top\hat\Sigma_\eta^{-1}\hat e\Rightarrow\chi^2_q\), since \(\Sigma_\eta\) is nonsingular. So \(\mathcal E\) holds for all \(\lambda\) at once, and for both ends, with probability \(\to1-\alpha\). Nonsingularity of \(\Sigma_\eta\) is what the earlier rank-\(q\) pseudoinverse argument lacked: there, a null-direction error of order \(o(n^{-1/2})\) is neither zero nor controlled. {[}alg; OLS; continuous mapping{]} ∎
\end{enumerate}

\emph{Remark (exact versus numerical reduction; added after review, round 2).} Case (ii-b) is claimed only for a documented \textbf{exact} reduction: the adding-up identity that removes column \(C\) (Remark P3.4.1), or coordinates fixed by the design. When the rank of \(\hat\Sigma\) comes from a numerical cutoff, the theorem applies only with \(q=\dim\beta\) (case (ii-a), assuming the true \(\Sigma\) is positive definite). Dropping candidates whose gradients load on a numerically null direction of \(\hat\Sigma\) is then a \textbf{safeguard}: it can only widen the interval (the minimum is over fewer terms), but it neither establishes the identities of (ii-b) nor calibrates a rank-\(q\) critical value. \emph{(An earlier response described that drop rule as case (ii-b); it is not.)}

\begin{enumerate}
\def\labelenumi{(\roman{enumi})}
\setcounter{enumi}{1}
\tightlist
\item
  needs no enumeration of vertices and is valid for any collection rule. It is conservative, increasingly so as \(q\) grows (many districts). (i) is sharper, but needs a fixed superset, and enumerating all dual vertices is feasible only in small problems. Checked (\path{verify_phase3.py}): weak duality on 500 random dual-feasible points. In a small LP, the projection bound built from the plug-in dual alone covered in 100\% of 3000 replications, against 94.5\% for a one-candidate \(z\)-bound.
\end{enumerate}

\phantomsection\label{pE-res-conjecture-P3.5-b3}\textbf{Conjecture P3.5(b3) (the implemented CLR procedure).} The package collects \(\hat\Lambda\) as the optimal dual vertices at \(\hat\beta\) and at \texttt{ndraw} Gaussian draws, and computes both \(k_p\) and \(\hat k\) over \(\hat\Lambda\). The conjecture is that it has asymptotic coverage \(\ge1-\alpha_U\) under (b1)'s conditions as \texttt{ndraw} \(\to\infty\).

\begin{itemize}
\tightlist
\item
  \emph{Proved reduction.} Suppose \(k_p\) is computed over \(\Lambda\), and \(\hat\Lambda\) contains every vertex selected by CLR's rule over \(\Lambda\) and a minimiser of \(\hat\theta_\mu+k_p\hat\sigma_\mu\) over \(\Lambda\). Then the threshold, the selected set, \(\hat k\) and \(\hat U\) all coincide with (b1), so (b1)'s coverage carries over. (The threshold is the same because its minimiser is in \(\hat\Lambda\). The selected set over \(\hat\Lambda\) is CLR's selected set intersected with \(\hat\Lambda\), which is CLR's selected set.)
\item
  \emph{Open.} (1) Recovery: vertices that are optimal only for rare Gaussian directions, degenerate optimal faces, and singular \(\Sigma\). (2) Computing \(k_p\) over \(\hat\Lambda\subseteq\Lambda\) can only lower it, so this error goes in the anti-conservative direction. (3) \(\lvert\Lambda\rvert\) grows with the number of districts, outside CLR's fixed-index setting.
\end{itemize}

Until this is settled, the CLR-type interval as implemented is an \emph{Algorithm} with simulation support (§8). The intervals of (a) and (b2) carry the formal guarantee.

\emph{Proof of (a).}

\begin{enumerate}
\def\labelenumi{\arabic{enumi}.}
\tightlist
\item
  Let \(\mathcal E\) be the event that \(s_j(\beta_0)\le s_j(\hat\beta)+\hat k\hat\sigma_{s_j}\) for all \(j\) and \(\lvert c(\beta_0)-c(\hat\beta)\rvert\le\hat k\hat\sigma_c\). Then \(P(\mathcal E)\to\ge1-\alpha_U-\alpha_L\). {[}OLS; linear maps of an asymptotically normal \(\hat\beta\), uniformly over the convex events of §0.5 (BR); Def of \(\hat k\); \(\hat\Sigma\) consistent. For M1/M2 this needs \(\dim\beta\) fixed, i.e.~R1{]}
\item
  On \(\mathcal E\), the true \(k_0\) satisfies \(s(\beta_0)+Ak_0\ge0\), hence \(s(\hat\beta)+\hat k\hat\sigma_s+Ak_0\ge0\), i.e.~\(k_0\in\hat\Theta^+\). {[}2b:Thm 2 / 2a:P1, the true value is feasible; \(\mathcal E\){]}
\item
  On \(\mathcal E\), \(\tilde B=c(\beta_0)+\gamma^\top k_0\) lies in \([c(\hat\beta)-\hat k\hat\sigma_c+\min_{\hat\Theta^+}\gamma^\top k,\ c(\hat\beta)+\hat k\hat\sigma_c+\max_{\hat\Theta^+}\gamma^\top k]\). {[}2; \(\mathcal E\); 2b:D2 with sample weights{]}
\item
  \(E=\sum_iw_ie_i\) with \(e_i\) independent given the design, mean 0, and \(\operatorname{Var}(e_i\mid V_i)\le\tfrac14\). So \(\operatorname{Var}(E\mid\text{design})\le S_1^2\). {[}P3-S1; CE; Pop{]}
\item
  \(\limsup_nP(\lvert E\rvert>z_{\alpha_E}S_1)\le\alpha_E\), provided \(\max_iw_i^2/\sum_iw_i^2\to0\). Let \(s_n=\operatorname{sd}(E\mid\text{design})\le S_1\) (step 4). It suffices that every subsequence has a further subsequence along which the \(\limsup\) is at most \(\alpha_E\). Take one along which \(s_n/S_1\to\rho\in[0,1]\).

  \begin{itemize}
  \tightlist
  \item
    If \(\rho=0\): by Chebyshev, \(P(\lvert E\rvert>z_{\alpha_E}S_1)\le s_n^2/(z_{\alpha_E}^2S_1^2)\to0\).
  \item
    If \(\rho>0\): the summands \(w_ie_i\) are independent with \(\lvert w_ie_i\rvert\le\lvert w_i\rvert\), and \(\max_i\lvert w_i\rvert/s_n\le2\max_i\lvert w_i\rvert/\big(\rho'(\textstyle\sum_iw_i^2)^{1/2}\big)\to0\) for any \(0<\rho'<\rho\) and \(n\) large. So Lindeberg's condition holds relative to \(s_n\), \(E/s_n\Rightarrow N(0,1)\), and \(P(\lvert E\rvert>z_{\alpha_E}S_1)\le P(\lvert E\rvert/s_n>z_{\alpha_E})\to\alpha_E\).
  \end{itemize}

  {[}Lyap; Chebyshev; 4{]} \emph{(Corrected after review, Batch B, Nov 2026: the earlier step applied a CLT under the weight condition alone, which fails when, e.g., only one summand has positive variance. The subsequence argument, Chebyshev where \(s_n/S_1\to0\) and Lindeberg where it is bounded away from 0, is the referee's; the conclusion and the allowance \(z_{\alpha_E}S_1\) are unchanged.)}
\item
  If \(B\notin\mathrm{CI}\), then either \(\tilde B\) lies outside the interval in 3, or \(\lvert E\rvert>z_{\alpha_E}S_1\). {[}\(B=\tilde B+E\); triangle inequality{]}
\item
  So \(P(B\notin\mathrm{CI})\le P(\mathcal E^c)+P(\lvert E\rvert>z_{\alpha_E}S_1)\to\le\alpha_U+\alpha_L+\alpha_E\). {[}union; 1, 3, 5, 6{]} ∎
\end{enumerate}

\emph{(b1)} replaces step 1 by CLR's coverage of each intersection bound, which needs the finite index set \(\Lambda\) (Lemma P3.1) and their Condition P {[}CLR{]}. Steps 4--7 are unchanged. \emph{(b2)} is proved above. The simulation-based collection of candidates is the subject of Conjecture P3.5(b3); it is not claimed as a theorem.

\textbf{Remarks.}

\begin{itemize}
\tightlist
\item
  \textbf{Better variance bound for the residual term.} \(\operatorname{Var}(e_{irc}\mid V)\le m(1-m)\) {[}BD{]}, with \(m=E[\beta_{irc}\mid V]\), which is bounded through the identified set. A data-dependent version needs its own coverage argument: Proposition P3.5(c) below.
\end{itemize}

\phantomsection\label{pE-res-proposition-P3.5-c}\textbf{Proposition P3.5(c) (variance-adaptive residual allowance; added after review, round 2, Nov 2026).} Fix a cell \((r,c)\) with weights \(w_i=w_{ir}\), put \(\pi_i=w_i^2/\sum_jw_j^2\) and \(\bar m=\sum_i\pi_im_{rc}(V_i)\), the \(\pi\)-weighted model-implied rate (a linear target in \((\beta,k)\)). By concavity of \(m(1-m)\),
\[s_n^2:=\operatorname{Var}(E\mid\text{design})\le\sum_iw_i^2m_i(1-m_i)\le\Big(\sum_iw_i^2\Big)\bar m(1-\bar m)=:\bar s^2 .\]
Let \(\mathrm{CI}_v\) be a confidence interval for \(\bar m\) at level \(1-\alpha_v\), and \(\hat S_{BD}=\big(\sum_iw_i^2\max_{m\in\mathrm{CI}_v\cap[0,1]}m(1-m)\big)^{1/2}\).

\begin{itemize}
\tightlist
\item
  \textbf{(c1) Valid version.} Suppose \(\mathrm{CI}_v\) is a P3.5(a) interval or a P3.5(b2)(ii) projection interval for \(\bar m\), computed on the \textbf{rung's own constraint system} (no added sign or CR rows), and suppose \(\max_iw_i^2/\bar s^2\to0\) (Lindeberg relative to \(\bar s\); it holds, for instance, when \(\bar m\) stays bounded away from 0 and 1). Replacing \(z_{\alpha_E}S_1\) by \(z_{\alpha_E}\hat S_{BD}\) in P3.5(a), or in the projection interval for \(B\), gives
  \[\liminf_nP(B_{rc}\in\mathrm{CI})\ge1-\alpha_{\rm est}-\alpha_E-\alpha_v ,\]
  with \(\alpha_{\rm est}\) the level spent on the estimation part. The Popoviciu version (\(S_1\), no \(\alpha_v\)) needs only \(\max_iw_i^2/\sum_iw_i^2\to0\) (step 5 of (a)).
\item
  \textbf{(c2) Adaptive version.} If \(\mathrm{CI}_v\) is the CLR interval (Conjecture P3.5(b3)), the same bound holds conditional on \(\mathrm{CI}_v\) covering \(\bar m\); the result inherits the algorithmic status of (b3).
\end{itemize}

\emph{Proof of (c1).}

\begin{enumerate}
\def\labelenumi{\arabic{enumi}.}
\tightlist
\item
  \(\bar m\) is a linear target under the rung (2a Prop. 7 with weights \(\pi\)), and the rung alone holds at the truth, so \(\mathcal E_v=\{\bar m\in\mathrm{CI}_v\}\) has \(\liminf P(\mathcal E_v)\ge1-\alpha_v\). {[}P3.5(a) or (b2)(ii){]}
\item
  On \(\mathcal E_v\), \(\hat S_{BD}\ge\bar s\ge s_n\). {[}concavity above; Def \(\hat S_{BD}\){]}
\item
  So \(\{\lvert E\rvert>z_{\alpha_E}\hat S_{BD}\}\subseteq\{\lvert E\rvert>z_{\alpha_E}\bar s\}\cup\mathcal E_v^c\), with \(\bar s\) deterministic given the design. {[}2{]}
\item
  \(\limsup P(\lvert E\rvert>z_{\alpha_E}\bar s)\le\alpha_E\): along a subsequence with \(s_n/\bar s\to0\) use Chebyshev; otherwise \(\max_i\lvert w_i\rvert/s_n\to0\) by the Lindeberg condition relative to \(\bar s\), and Lindeberg's CLT applies relative to \(s_n\), as in step 5 of (a) with \(\bar s\) in place of \(S_1\). {[}Chebyshev; Lyap{]}
\item
  Combine with the estimation event of (a) or (b2)(ii) by the union bound. {[}union; 1--4{]} ∎
\end{enumerate}

\emph{Why the condition matters.} If \(\mathrm{CI}_v\) fails, BD is not conservative: with true \(\bar m=.5\) a wrong interval inside \([0,.1]\) gives a variance bound at most \(.09\) against the true \(.25\). And if \(\bar s/S_1\to0\) (e.g.~\(\bar m\to0\)), Lindeberg relative to \(\bar s\) can fail, so step 4 needs its condition. Check: \path{verify_phase3.py}, ``P3.5(c)''. \emph{(The referee asked that the provably valid projection interval for \(B\) be separated from the variance-adaptive one. In the package (0.3.x), the projection column for \(B\) uses (c1) with \(\mathrm{CI}_v\) the projection interval on the rung alone, and prints a note when \(\mathrm{CI}_v\) comes near 0 or 1; the CLR \texttt{ci\_actual} uses (c2) and is labelled algorithmic. \(\tilde B\) is unaffected.)}

\begin{itemize}
\tightlist
\item
  \textbf{Optional assumption for a much smaller residual term.} If the poll-level deviations of different groups are not negatively correlated, then \(\operatorname{Var}(X_{ir}e_{irc})\le\operatorname{Var}(T_{ic}\mid V_i)\), which is estimable. In simulation this would remove most of the residual allowance. It is an added assumption and belongs on the menu, not in the default.
\item
  \textbf{JKST's \(\mathrm{CI}_x\)} is version (b) with a single candidate per side (tie-breaking), no selection step, and \(S_L=S_1+\hat\sigma_{\lambda^*}\), which combines the two error sources by adding standard deviations instead of by the union bound.
\end{itemize}

\phantomsection\label{pE-res-proposition-P3.6}\textbf{Proposition P3.6 (tests).}

\begin{enumerate}
\def\labelenumi{(\alph{enumi})}
\tightlist
\item
  \textbf{Emptiness test.} Reject the model (rung, foundations) at level \(\alpha_U+\alpha_L\) if \(\hat\Theta^+\) is empty.
\item
  \textbf{Test of M1's testable part.} Run a Wald/F test of the M1 regression against M2 (the overidentifying restrictions of 2b D1(b)), with \((D^\star-1)(C-1)R(R-1)/2\) restrictions. In 2×2: equal curvature across districts.
\item
  \textbf{Test of M0 against M1 (pooling).} Conditional on the design, suppose each district's observed compositions contain \(R\) affinely independent points. Under M1, M0's mean model holds for the observed units iff all district shifts vanish, \(\zeta_{c,d}=0\). So pooling is tested by a Wald/F test of \(\zeta=0\) in the M1 regression, with \((D^\star-1)(C-1)R\) restrictions. \emph{(Added after review. 2b Prop. D4(b) gives a weaker condition, \(E[Z_D\mid X,N]\) affine in \(X\). That condition averages over the district labels, so it describes the superpopulation pooled over re-drawn districts, which §0.4 does not target. Conditional on the observed labels, only \(\zeta=0\) makes M0 correct.)} (c) tests mean shifts only. It has no power against within-district residual dependence without mean shifts, which violates P3-S1a (§0.4). \emph{(Sentence added after review, round 2.)}
  \emph{Proof of (a).}
\end{enumerate}

\begin{enumerate}
\def\labelenumi{\arabic{enumi}.}
\tightlist
\item
  Under the model, \(k_0\in\hat\Theta^+\) on \(\mathcal E\). {[}P3.5(a) step 2{]}
\item
  So \(P(\hat\Theta^+=\emptyset)\le P(\mathcal E^c)\to\le\alpha_U+\alpha_L\). {[}1{]}
\end{enumerate}

\begin{enumerate}
\def\labelenumi{(\alph{enumi})}
\setcounter{enumi}{1}
\tightlist
\item
  is the standard Wald test {[}OLS{]}. Its \(\chi^2\) calibration is established for fixed dimension (R1). Under R2/R3 the number of restrictions grows with \(D^\star\). The F-form with HC3, or a wild bootstrap, is then a practical choice with simulation support only.
\end{enumerate}

\textbf{Calibration in the package (0.3.0; added after the Batch B review and the WS6 Monte Carlo, Nov 2026).} In moderate dimension the HC3 Wald statistic against \(\chi^2\) over-rejects: for the pooling test (c) in design S0 of §8.3 (38 restrictions), 0.078 at nominal 5\%, and 0.072 for the F form. The package therefore calibrates both (b) and (c) by a \textbf{restricted wild bootstrap}:

\begin{itemize}
\tightlist
\item
  fit the null model (M0 for (c); M1 for (b));
\item
  draw \(Y^*=\hat Y_0+\hat e_0v\) with Rademacher \(v_i\), one per unit, shared across the \(C-1\) equations;
\item
  refit the larger model and recompute the same HC Wald statistic; \(p=(1+\#\{W^*\ge W\})/(B+1)\), with \(B=399\) by default.
\end{itemize}

The \(\chi^2\) and F p-values are reported as well. Rejection rates at nominal 5\% in design S0 (MC standard errors in parentheses):

\begin{longtable}[]{@{}
  >{\raggedright\arraybackslash\hspace{0pt}}p{(\columnwidth - 10\tabcolsep) * \real{0.1667}}
  >{\raggedright\arraybackslash\hspace{0pt}}p{(\columnwidth - 10\tabcolsep) * \real{0.1667}}
  >{\raggedright\arraybackslash\hspace{0pt}}p{(\columnwidth - 10\tabcolsep) * \real{0.1667}}
  >{\raggedright\arraybackslash\hspace{0pt}}p{(\columnwidth - 10\tabcolsep) * \real{0.1667}}
  >{\raggedright\arraybackslash\hspace{0pt}}p{(\columnwidth - 10\tabcolsep) * \real{0.1667}}
  >{\raggedright\arraybackslash\hspace{0pt}}p{(\columnwidth - 10\tabcolsep) * \real{0.1667}}@{}}
\toprule\noalign{}
\begin{minipage}[b]{\linewidth}\raggedright\hspace{0pt}Test
\end{minipage} & \begin{minipage}[b]{\linewidth}\raggedright\hspace{0pt}Restrictions
\end{minipage} & \begin{minipage}[b]{\linewidth}\raggedright\hspace{0pt}Replications
\end{minipage} & \begin{minipage}[b]{\linewidth}\raggedright\hspace{0pt}\(\chi^2\)
\end{minipage} & \begin{minipage}[b]{\linewidth}\raggedright\hspace{0pt}F
\end{minipage} & \begin{minipage}[b]{\linewidth}\raggedright\hspace{0pt}Wild bootstrap
\end{minipage} \\
\midrule\noalign{}
\endhead
\bottomrule\noalign{}
\endlastfoot
(c) pooling, M0 vs M1 & 38 & 2000 (\(B=399\)) & 0.078 (0.006) & 0.072 (0.006) & \textbf{0.043} (0.005) \\
(b) M1 vs M2 & 19 & 1000 (\(B=199\)) & 0.057 (0.007) & 0.055 (0.007) & \textbf{0.054} (0.007) \\
\end{longtable}

Scripts: \path{benchmarks/pooling_test_size.R}, \path{benchmarks/m1_test_size.R} (outputs in \path{benchmarks/results/}). The bootstrap is a practical calibration supported by simulation; its validity as the number of restrictions grows is not proved.

\emph{Proof of (c).}

\begin{enumerate}
\def\labelenumi{\arabic{enumi}.}
\tightlist
\item
  Under M1 and P3-S1a, \(E[\beta_i\mid\text{design}]=M_H(X_i)+Z_{D_i}\), with \(Z_d=[\zeta_{c,d,r}]_{rc}\). {[}Def M1; 2b §0.1{]}
\item
  M0 holds for the observed units iff there is \(H'\) with \(M_{H'}(X_i)=M_H(X_i)+Z_{D_i}\) for all \(i\), i.e.~\(G_cX_i=\zeta_{c,D_i}\) with \(G=H'-H\). {[}Def M0{]}
\item
  Within district \(d\), \((G_c-\zeta_{c,d}\mathbf 1^\top)X_i=0\) for every \(i\) in \(d\) (using \(\mathbf 1^\top X_i=1\)). The \(X_i\) include \(R\) affinely independent points, which are linearly independent because they lie in \(\{\mathbf 1^\top x=1\}\). So \(G_c=\zeta_{c,d}\mathbf 1^\top\). {[}alg{]}
\item
  This holds for every \(d\), and \(\zeta_{c,1}=0\), so \(G_c=0\) and every \(\zeta_{c,d}=0\). Conversely, \(\zeta=0\) gives M0 with \(H'=H\). {[}3; Def M1{]}
\item
  The \(\zeta\) are coefficients of the M1 regression, identified by 2b Lemma D1(d). The Wald test of \(\zeta=0\) has \((D^\star-1)(C-1)R\) restrictions for the \(C-1\) free columns. {[}D1(d); OLS{]} ∎
\end{enumerate}

\hypertarget{pE-a-finite-sample-interval-valid-in-every-regime}{%
\subsubsection{5.1 A finite-sample interval, valid in every regime}\label{pE-a-finite-sample-interval-valid-in-every-regime}}

\phantomsection\label{pE-res-theorem-P3.8}\textbf{Theorem P3.8 (finite-sample CI and emptiness test).} Fix a rung and condition on the design. Assume P3-S1a, that the rung's mean model holds, and that its regression design \(Z\) has full column rank. Fit the \(C-1\) free columns by OLS, \(\hat\beta_c=(Z^\top Z)^{-1}Z^\top T_{\cdot c}\).

\emph{Hoeffding radius of a linear functional.} Every linear functional \(f(\beta)=f_0+\sum_{c<C}b_c^\top\beta_c\) satisfies \(f(\hat\beta)=f_0+\sum_i\sum_{c<C}w_{ic}(f)T_{ic}\), with \(w_{ic}(f)=b_c^\top(Z^\top Z)^{-1}z_i\). Put \(w_{iC}(f):=0\), \(\rho_i(f):=\max_cw_{ic}(f)-\min_cw_{ic}(f)\), and
\[h(f;\delta)=\Big(\tfrac12\log(1/\delta)\sum_i\rho_i(f)^2\Big)^{1/2}.\]

\emph{Interval.} Let \(J\) be the number of constraint rows and \(\delta=\alpha_s/(J+2)\). Put \(h_j=h(s_j;\delta)\), \(h_c=h(c;\delta)\), \(t_E=\big(\tfrac12\log(2/\alpha_E)\sum_iw_{ir}^2\big)^{1/2}\), and
\[\hat\Theta^{\rm fin}=\{k:s_j(\hat\beta)+h_j+a_j^\top k\ge0\ \ \forall j\},\]
\[\mathrm{CI}^{\rm fin}=\Big[c(\hat\beta)-h_c+\min_{\hat\Theta^{\rm fin}}\gamma^\top k-t_E,\ \ c(\hat\beta)+h_c+\max_{\hat\Theta^{\rm fin}}\gamma^\top k+t_E\Big].\]
Then, \textbf{for every \(n\)} and in every regime:

\begin{enumerate}
\def\labelenumi{(\alph{enumi})}
\tightlist
\item
  \(P(B_{rc}\in\mathrm{CI}^{\rm fin}\mid\text{design})\ge1-\alpha_s-\alpha_E\);
\item
  without \(t_E\), \(P(\tilde B_{rc}\in\mathrm{CI}^{\rm fin}\mid\text{design})\ge1-\alpha_s\);
\item
  rejecting the rung when \(\hat\Theta^{\rm fin}=\emptyset\) is a test of exact level at most \(\alpha_s\);
\item
  intersecting with Duncan--Davis keeps (a), since \(B\in DD\) always.
\end{enumerate}

\emph{Proof.}

\begin{enumerate}
\def\labelenumi{\arabic{enumi}.}
\tightlist
\item
  \(T_i\in\Delta^{C-1}\), so \(Y_i(f):=\sum_cw_{ic}(f)T_{ic}\) is a convex combination of the \(w_{ic}(f)\). It lies in an interval of length \(\rho_i(f)\). {[}S1; Def{]}
\item
  Under the mean model and full rank, OLS is conditionally unbiased: \(E[f(\hat\beta)\mid\text{design}]=f(\beta_0)\). {[}linearity of OLS; \(E[T_{ic}\mid\text{design}]=z_i^\top\beta_{c,0}\){]}
\item
  The \(Y_i(f)\) are independent given the design, so \(P\big(f(\beta_0)-f(\hat\beta)\ge h(f;\delta)\big)\le\delta\), and the same for the other side. {[}P3-S1a; Hoeff; 1--2{]}
\item
  Let \(\mathcal E\) be the event that \(s_j(\beta_0)\le s_j(\hat\beta)+h_j\) for all \(j\) and \(\lvert c(\beta_0)-c(\hat\beta)\rvert\le h_c\). It is an intersection of \(J+2\) one-sided events, so \(P(\mathcal E^c)\le(J+2)\delta=\alpha_s\). {[}3; union{]}
\item
  On \(\mathcal E\), the true \(k_0\) satisfies \(s(\beta_0)+Ak_0\ge0\), hence \(k_0\in\hat\Theta^{\rm fin}\), and \(\tilde B=c(\beta_0)+\gamma^\top k_0\) lies in the interval without \(t_E\). That proves (b), and (c), since \(\hat\Theta^{\rm fin}\ne\emptyset\) on \(\mathcal E\). {[}P3.5(a) steps 2--3{]}
\item
  \(E=\sum_iw_{ir}e_{irc}\) with independent \(e_{irc}\in[-m_{rc}(V_i),1-m_{rc}(V_i)]\), an interval of length 1. So \(P(\lvert E\rvert\ge t_E)\le2\exp(-2t_E^2/\sum_iw_{ir}^2)=\alpha_E\). {[}P3-S1a; Hoeff{]}
\item
  \(B=\tilde B+E\). A union bound over 4 and 6 gives (a). (d) holds because \(B\in DD\) surely. {[}union; 2b:D0{]} ∎
\end{enumerate}

\emph{Remarks.}

\begin{itemize}
\tightlist
\item
  \textbf{What it needs, and what it does not.}

  \begin{itemize}
  \tightlist
  \item
    Needs: independence given the design (P3-S1a), a correctly specified mean model, and a full-rank design.
  \item
    Does not need: Gaussian approximation, a fixed dimension, Slater, or any regime.
  \end{itemize}

  It is therefore the only formal guarantee for M1 with many small districts (R2), where the asymptotic results are not established.
\item
  \textbf{Width: valid but currently uninformative.} In a 2×2 test design with linear means, Beta-distributed rates and composition ranges of 0.2--0.45 per district, the interval clipped to \([0,1]\) had median width 1.0 in every configuration tried:

  \begin{itemize}
  \tightlist
  \item
    pooled, with 900 and with 5,000 polls;
  \item
    5 districts × 200 polls;
  \item
    10 × 93;
  \item
    30 × 8.
  \end{itemize}

  Two costs combine:

  \begin{enumerate}
  \def\labelenumi{\arabic{enumi}.}
  \tightlist
  \item
    Hoeffding uses the worst-case variance \(\tfrac14\) for each \(T_{ic}\), where the true variance was about 0.01.
  \item
    The separate union bound over all \(J\) constraints and the constant inflates each radius, and each radius is then divided by the distance from the vertex to a pure composition.
  \end{enumerate}

  An \emph{oracle} Gaussian union bound with the true variances still gave widths of 0.76--0.87, and 1.0 at 30 × 8. So most of the cost is the constraint-by-constraint union bound, which the conservative asymptotic interval P3.5(a) shares. Theorem P3.8 settles validity in R2, but it does not yet give useful intervals there. Informative finite-sample inference for R2 is open (§10).
\end{itemize}

Checked (\path{verify_phase3.py}): in the 30 × 8 design, coverage of \(\tilde B\) and \(B\) was 100\% in 300 replications each, and the interval was never empty. In the 1000-replication Monte Carlo (§8.2) it covered in every replication of five M1 designs (10 × 93, 93 × 30, and 30, 93 and 200 districts × 10 polls), with clipped width 1.000 in all of them.

\begin{itemize}
\tightlist
\item
  \textbf{Possible refinements, not proved here.}

  \begin{itemize}
  \tightlist
  \item
    An empirical-Bernstein bound (Maurer \& Pontil 2009) would replace the ranges \(\rho_i\) by estimated variances.
  \item
    A Bhatia--Davis version of \(t_E\), as in §5.
  \item
    Splitting \(\alpha_s\) unevenly across constraints, e.g.~giving less to constraints that cannot bind.
  \end{itemize}
\end{itemize}

\begin{center}\rule{0.5\linewidth}{0.5pt}\end{center}

\hypertarget{pE-many-districts-with-few-units-regime-r2}{%
\subsection{6. Many districts with few units (regime R2)}\label{pE-many-districts-with-few-units-regime-r2}}

\phantomsection\label{pE-res-proposition-P3.7}\textbf{Proposition P3.7.}

\begin{enumerate}
\def\labelenumi{(\alph{enumi})}
\tightlist
\item
  \emph{(Dropping constraints is conservative.)} For any subset \(\mathcal D'\) of districts, \(\Theta^{M1}_O(\mathcal D')\supseteq\Theta^{M1}_O\).
\item
  \emph{(Design-based selection preserves conditional validity.)} Suppose \(\mathcal D'\) is chosen as a function of the design \(\{(N_i,X_i,D_i)\}\) only. Then any procedure that has (finite-sample or asymptotic) coverage when applied with a fixed set of districts keeps that coverage conditionally on the selected \(\mathcal D'\). In particular, Theorem P3.8 holds for any design-selected \(\mathcal D'\). The asymptotic results of §§3--5 transfer only if their own regime conditions (§0.5) hold for the retained districts. Selection does not by itself supply them: retained districts with bounded \(n_d\) are still in R2.
  \emph{Proof.}
\end{enumerate}

\begin{enumerate}
\def\labelenumi{\arabic{enumi}.}
\tightlist
\item
  Fewer constraints give a weakly larger feasible set. {[}Def; 2b:Thm 2(b){]}
\item
  Conditioning on the design fixes \(\mathcal D'\). A procedure's conditional coverage given the design, for the fixed set \(\mathcal D'\), is therefore its coverage given the design and \(\mathcal D'\). {[}P3-S2{]} ∎
\end{enumerate}

\emph{(Revised after review: the earlier statement said every result of §§3--5 holds conditionally on \(\mathcal D'\), which overstated what selection achieves.)}

\textbf{What happens under R2.}

\begin{itemize}
\tightlist
\item
  Each district's shifts \(\hat\zeta_d\) come from about \(n_d\) units, and their sampling error does not shrink.
\item
  The number of candidates grows with \(D^\star\).
\item
  The plug-in is badly too narrow. In the Monte Carlo (§8.2), with 10 polls per district, it covers 31\%, 4\% and 0.5\% at 30, 93 and 200 districts, and is empty 24\%, 46\% and 72\% of the time. This is the Theorem P3.4(ii) bias, compounded across many districts.
\item
  In the Monte Carlo (§8.2), the CLR interval covered 0.995--1 in every R2 design, up to 200 districts. This is a \emph{simulation finding} (§0.5), not a theorem. CLR inflates each candidate in proportion to its own standard error, so noisy districts rarely bind, which is precision weighting built into the method.
\item
  The formal guarantee in R2 is Theorem P3.8, at the cost of width: in §8.2 it is \([0,1]\) in every design. The projection interval is also \([0,1]\) there, because \(q\) grows with \(D^\star\).
\item
  With 10 units per district, CLR's interval is about as wide as Duncan--Davis at 30--93 districts, so the identification gain from district dummies (2b) is consumed by sampling noise. It narrows as districts are added, and is below DD at 200 districts (0.42 for \(\tilde B\) and 0.49 for \(B\), against 0.57).
\end{itemize}

Practical options. Each keeps the identified set valid (P3.7(a)). Their inference is formally covered by Theorem P3.8, and asymptotically only under §0.5's conditions:

\begin{itemize}
\item
  \begin{enumerate}
  \def\labelenumi{(\roman{enumi})}
  \tightlist
  \item
    Keep the constraints of districts with enough units or composition spread (P3.7).
  \end{enumerate}
\item
  \begin{enumerate}
  \def\labelenumi{(\roman{enumi})}
  \setcounter{enumi}{1}
  \tightlist
  \item
    Use M1(\(\delta\)), \(\delta>0\), which bounds how far any one district can pull the shared parameter.
  \end{enumerate}
\item
  \begin{enumerate}
  \def\labelenumi{(\roman{enumi})}
  \setcounter{enumi}{2}
  \tightlist
  \item
    Pool districts into larger groups.
  \end{enumerate}
\item
  \begin{enumerate}
  \def\labelenumi{(\roman{enumi})}
  \setcounter{enumi}{3}
  \tightlist
  \item
    Target \(\tilde B\) to drop the residual term.
  \end{enumerate}
\end{itemize}

A hierarchical restriction on the district shifts (they are draws from a distribution) would change the model and is left open.

\begin{center}\rule{0.5\linewidth}{0.5pt}\end{center}

\hypertarget{pE-sensitivity-analysis}{%
\subsection{7. Sensitivity analysis}\label{pE-sensitivity-analysis}}

\phantomsection\label{pE-res-proposition-S1}\textbf{Proposition S1 (bounded nonlinearity).} Replace (S2) by \(E[\beta\mid V]=M_H(X)+\Delta(X)\), with \(\Delta\) in the undetectable class of 2a Prop. 9 (\(x^\top\Delta(x)=0\), \(\Delta(x)\mathbf 1=0\)) and \(\lVert\Delta(x)\rVert_\infty\le\eta\). Then \(Q\) is still identified as before, and every value of \(\tilde B_{rc}\) consistent with the data lies in
\[\Big[\min_{\Theta_O(\eta)}\varphi-\eta,\ \max_{\Theta_O(\eta)}\varphi+\eta\Big],\qquad\Theta_O(\eta)=\{k:s+Ak\ge-\eta\mathbf 1\}.\]
For a general \(\varphi=\sum a_{rc}B_{rc}\), widen by \(\eta\sum\lvert a_{rc}\rvert\) instead. This is an \textbf{outer} bound under bounded deviations: it is valid, but it is not claimed to be the sharp set of the relaxed model.

\emph{Proof.}

\begin{enumerate}
\def\labelenumi{\arabic{enumi}.}
\tightlist
\item
  \(E[T\mid X]=X^\top Q X+X^\top\Delta(X)=X^\top QX\), so \(Q\) is unchanged. {[}2a:L1(a); \(x^\top\Delta=0\){]}
\item
  \(M_H(x)+\Delta(x)\in\mathfrak B\) implies \(M_H(x)\ge-\Delta(x)\ge-\eta\) elementwise on \(A\). \(M_H\) is linear, so this is equivalent to the same inequality at the vertices: \(s+Ak\ge-\eta\). {[}S3; 2a:L3(a){]}
\item
  \(\tilde B_{rc}=\sum_iw_i[M_H(X_i)+\Delta(X_i)]_{rc}=\varphi(k)+\sum_iw_i\Delta_{rc}(X_i)\), and the last term lies in \([-\eta,\eta]\) because the weights sum to 1. {[}2b:D2; Def{]} ∎
\end{enumerate}

Checked: a nonlinear undetectable violation with \(\eta=0.02\) stays inside the bound, and the bounds are nested in \(\eta\).

\textbf{Breakdown value.} For a conclusion \(\mathcal C\), e.g.~``loyalty \textgreater{} 0.30'', define \(\eta^*(\mathcal C)=\sup\{\eta:\mathcal C\text{ holds for every value in the }\eta\text{-bound}\}\). The bounds are nested in \(\eta\), so the conclusion holds for all \(\eta<\eta^*\) and fails above it, and \(\eta^*\) can be found by bisection. In the 3×3 example, \(\eta^*(\text{loyalty}>0.30)=0.060\): the conclusion survives any undetectable nonlinearity of up to 6 percentage points. Apply the CI version of the bound to account for sampling error.

\textbf{Other dials, same logic} (nested sets, bisection for breakdown):

\begin{itemize}
\tightlist
\item
  \(\delta\) for district heterogeneity (2b Prop. D6);
\item
  \(\delta\) for change between periods (menu item 19);
\item
  population turnover \(\epsilon\) (A1). Turnover requires relaxing the Duncan--Davis sets and the regression identity together. It is left as an open item, with interim advice to add explicit entry/exit rows.
\end{itemize}

\begin{center}\rule{0.5\linewidth}{0.5pt}\end{center}

\hypertarget{pE-monte-carlo}{%
\subsection{8. Monte Carlo}\label{pE-monte-carlo}}

\emph{Rerun for review item R-37 (Nov 2026). It replaces the earlier 200-replication table.}

\path{phase3_montecarlo.py}, \textbf{1000 replications per design} (\path{phase3_montecarlo_output.txt}, summary in \path{phase3_montecarlo_summary.json}). Unit rates are Beta/Dirichlet draws whose conditional means are exactly linear in composition given the district. So M1 always holds, and M0 holds only when the district shifts are zero. Every interval is nominal 95\%:

\begin{itemize}
\tightlist
\item
  \textbf{plug-in}: the estimate of the identified set (not a CI);
\item
  \textbf{conservative} (Thm P3.5(a));
\item
  \textbf{CLR} as implemented (Conjecture P3.5(b3)), for \(B\) with the Bhatia--Davis allowance (the package default) and for \(\tilde B\);
\item
  \textbf{projection} (Prop. P3.5(b2)(ii)): \(k=\sqrt{\chi^2_{q,1-\alpha}}\) with \(q=\operatorname{rank}\hat\Sigma\), applied to the same candidates as CLR. In every 2×2 design here \(\hat\Sigma\) has full rank, so \(q=\dim\beta\) and this is case (ii-a) of the corrected proposition. (Since 0.3.0 the package also drops candidates whose gradients use a numerical null direction of \(\hat V\). \emph{Corrected after review, round 2:} that is a safeguard, not case (ii-b); see the remark after P3.5(b2).) For \(B\): estimation part at level 0.04, plus the Bhatia--Davis allowance with \(\bar m\) bounded by the \textbf{projection} interval at level \(\alpha_V=0.005\) on the rung alone (P3.5(c1); column \texttt{proj\_B}). The earlier runs took \(\bar m\) from a CLR interval; that version is kept as \texttt{proj\_B\_clrbd}, labelled \textbf{not covered} (P3.5(c2));
\item
  \textbf{finite-sample} (Thm P3.8), in the 2×2 M1 designs.
\end{itemize}

Conventions:

\begin{itemize}
\tightlist
\item
  \textbf{Coverage} is the share of 1000 replications whose interval contains the truth. The Monte Carlo standard error, \(\sqrt{p(1-p)/1000}\), is in parentheses. A coverage of \textbf{1} means all 1000 replications covered; the exact one-sided 95\% lower confidence bound is then \(0.05^{1/1000}=0.997\).
\item
  An empty interval (plug-in or conservative set empty, or CLR/projection ends crossing) counts as not covering.
\item
  \textbf{Width} is the median after clipping to \([0,1]\). The Duncan--Davis (DD) width is shown for reference.
\item
  Each cell reads ``coverage / width''.
\end{itemize}

\textbf{Cross-check against the package.} \path{phase3_montecarlo_crosscheck.py} runs one replication of designs A (10 × 93) and S2 through \texttt{eiBounds()} (\path{phase3_montecarlo_crosscheck_output.txt}). The two agree exactly on the plug-in, the projection interval for \(\tilde B\) and the pooling-test statistic. The CLR ends differ by up to 0.01, because the script uses every constraint as a candidate (the full dual-vertex set in 2×2) while the package collects candidates by simulation. The package also intersects \texttt{ci\_actual}/\texttt{proj\_actual} with DD.

\textbf{Which figures correspond to proved results.}

\begin{itemize}
\item
  Conservative and projection intervals: Thm P3.5(a) and Prop. P3.5(b2)(ii) are asymptotic theorems for a fixed-dimensional regression: M0 in every design here, and M1 with few, large districts (A 10 × 93; S 20 × 50). In the R2 rows (93 × 30, and 10 polls per district) \(q\) grows with \(D^\star\), and their coverage there is a simulation finding (§0.5).
\item
  Finite-sample interval: Thm P3.8 holds in every row.
\item
  CLR: Conjecture P3.5(b3) everywhere. All its figures are simulation findings.
\item
  The projection interval for \(B\) (\texttt{proj\_B}) uses the Bhatia--Davis allowance with \(\bar m\) bounded by a projection interval on the rung alone, so it is covered by P3.5(b2)(ii) and P3.5(c1), subject to (c1)'s Lindeberg condition. \emph{(Corrected after review, round 2, Nov 2026: this bullet said the variance bound came from a CLR interval ``as in the package''. Since 0.3.0 the package takes it from the projection interval, and the §8 column has been rerun to match (table below). The CLR-based version, \texttt{proj\_B\_clrbd}, is not covered by any theorem.)}
\item
  \textbf{Rerun (round 2, Nov 2026)} of the projection column for \(B\) with the P3.5(c1) bound, same seeds (no other column uses the changed code, and every other column reproduces the earlier run exactly; \path{phase3_montecarlo_output.txt}):

  \begin{longtable}[]{@{}
    >{\raggedright\arraybackslash\hspace{0pt}}p{(\columnwidth - 4\tabcolsep) * \real{0.3333}}
    >{\raggedright\arraybackslash\hspace{0pt}}p{(\columnwidth - 4\tabcolsep) * \real{0.3333}}
    >{\raggedright\arraybackslash\hspace{0pt}}p{(\columnwidth - 4\tabcolsep) * \real{0.3333}}@{}}
  \toprule\noalign{}
  \begin{minipage}[b]{\linewidth}\raggedright\hspace{0pt}  Design
  \end{minipage} & \begin{minipage}[b]{\linewidth}\raggedright\hspace{0pt}  \texttt{proj\_B} (P3.5(c1), \(\bar m\) from the projection interval): coverage / median width
  \end{minipage} & \begin{minipage}[b]{\linewidth}\raggedright\hspace{0pt}  \texttt{proj\_B\_clrbd} (\(\bar m\) from CLR; earlier column, \textbf{not covered})
  \end{minipage} \\
  \midrule\noalign{}
  \endhead
  \bottomrule\noalign{}
  \endlastfoot
  A 10 × 93 (M1) & 1.000 / 0.662 & 1.000 / 0.662 \\
  A 93 × 30 (M1) & 1.000 / 0.964 & 1.000 / 0.964 \\
  R2 30 × 10, 93 × 10, 200 × 10 (M1) & 1.000 / 1.000 each & 1.000 / 1.000 each \\
  C boundary (M0) & 1.000 / 0.059 & 1.000 / 0.059 \\
  E Slater (M0) & 1.000 / 0.251 & 1.000 / 0.249 \\
  S0, S1 (M0; M1) & 1.000 / 0.436, 0.444; 1.000 / 0.808, 0.773 & identical \\
  S2 non-affine (M0, misspecified) & 0.206 (SE 0.013) / 0.038, empty 0.478 & identical \\
  S2 (M1) & 1.000 / 0.611 & 1.000 / 0.611 \\
  B 3×3, four targets (M0) & 1.000 / 0.679, 0.445, 0.378, 0.661 & 1.000 / 0.679, 0.445, 0.377, 0.661 \\
  \end{longtable}

  Coverage is unchanged in every design, and widths change by at most 0.002 (design E and one design-B target). In these designs the projection interval for \(\bar m\) is only slightly wider than CLR's, and the Bhatia--Davis maximum \(\max m(1-m)\) over it is flat near the plausible rates. So the provably valid version costs essentially nothing here. The S2/M0 row is misspecification, not a failure of the theorem.
\end{itemize}

\textbf{Rank of \(\hat\Sigma\) in the projection interval} (\path{phase3_montecarlo_rank.py}, \path{phase3_montecarlo_rank_output.txt}; same seeds, 1000 replications per design). These runs use \(q=\operatorname{rank}\hat\Sigma\), which equals \(\dim\beta\) in every 2×2 design and the exact adding-up reduction in design B (below); after the round-2 review the theorem is claimed with \(q=\dim\beta\) or a documented exact reduction, and a numerical cutoff is only a safeguard (remark after P3.5(b2)). Since the Batch B correction P3.5(b2)(ii) requires either a positive definite \(\Sigma\) (ii-a) or a reduced parameterisation with exactly zero error in the null directions (ii-b). If \(\hat\Sigma\) were singular while the estimation error had a component in its null space, those directions would not be controlled. The numerical rank (eigenvalues above \(10^{-10}\) times the largest) equals \(\dim\beta\) in every replication of every 2×2 design, from \(\dim\beta=3\) to 401.

\begin{itemize}
\tightlist
\item
  Conditioning: the smallest-to-largest eigenvalue ratio is 0.005--0.04 for the pooled fits, and falls with the number of districts, to a minimum of \(2.6\times10^{-8}\) at 200 × 10. That is still far above the rank threshold.
\item
  Design B fits all three columns (\(\dim\beta=18\)) and has \(q=12\) in every replication. The six null directions are the adding-up constraint: summed across columns, the coefficients are the regression of 1 on the design, which has no error. So the estimation error lies in the range of \(\hat\Sigma\), and nothing is left uncontrolled.
\item
  On the \(C-1\) free columns (Remark P3.4.1, as the package fits), \(\hat\Sigma\) has full rank 12 in every replication.
\item
  No replication had a numerical rank drop, so there is no rank-drop subsample whose coverage could be compared.
\end{itemize}

\hypertarget{pE-model-holds-m1-or-m0-without-shifts}{%
\subsubsection{8.1 Model holds (M1, or M0 without shifts)}\label{pE-model-holds-m1-or-m0-without-shifts}}

\begin{longtable}[]{@{}
  >{\raggedright\arraybackslash\hspace{0pt}}p{(\columnwidth - 18\tabcolsep) * \real{0.1000}}
  >{\raggedright\arraybackslash\hspace{0pt}}p{(\columnwidth - 18\tabcolsep) * \real{0.1000}}
  >{\raggedright\arraybackslash\hspace{0pt}}p{(\columnwidth - 18\tabcolsep) * \real{0.1000}}
  >{\raggedright\arraybackslash\hspace{0pt}}p{(\columnwidth - 18\tabcolsep) * \real{0.1000}}
  >{\raggedright\arraybackslash\hspace{0pt}}p{(\columnwidth - 18\tabcolsep) * \real{0.1000}}
  >{\raggedright\arraybackslash\hspace{0pt}}p{(\columnwidth - 18\tabcolsep) * \real{0.1000}}
  >{\raggedright\arraybackslash\hspace{0pt}}p{(\columnwidth - 18\tabcolsep) * \real{0.1000}}
  >{\raggedright\arraybackslash\hspace{0pt}}p{(\columnwidth - 18\tabcolsep) * \real{0.1000}}
  >{\raggedright\arraybackslash\hspace{0pt}}p{(\columnwidth - 18\tabcolsep) * \real{0.1000}}
  >{\raggedright\arraybackslash\hspace{0pt}}p{(\columnwidth - 18\tabcolsep) * \real{0.1000}}@{}}
\toprule\noalign{}
\begin{minipage}[b]{\linewidth}\raggedright\hspace{0pt}Design
\end{minipage} & \begin{minipage}[b]{\linewidth}\raggedright\hspace{0pt}\(q\)
\end{minipage} & \begin{minipage}[b]{\linewidth}\raggedright\hspace{0pt}DD
\end{minipage} & \begin{minipage}[b]{\linewidth}\raggedright\hspace{0pt}Plug-in: cover / empty / width
\end{minipage} & \begin{minipage}[b]{\linewidth}\raggedright\hspace{0pt}Conservative, \(B\)
\end{minipage} & \begin{minipage}[b]{\linewidth}\raggedright\hspace{0pt}CLR, \(B\)
\end{minipage} & \begin{minipage}[b]{\linewidth}\raggedright\hspace{0pt}Projection, \(B\)
\end{minipage} & \begin{minipage}[b]{\linewidth}\raggedright\hspace{0pt}Conservative, \(\tilde B\)
\end{minipage} & \begin{minipage}[b]{\linewidth}\raggedright\hspace{0pt}CLR, \(\tilde B\)
\end{minipage} & \begin{minipage}[b]{\linewidth}\raggedright\hspace{0pt}Projection, \(\tilde B\)
\end{minipage} \\
\midrule\noalign{}
\endhead
\bottomrule\noalign{}
\endlastfoot
A: 2×2, M1, 10 districts × 93 & 21 & 0.574 & 0.982 (0.004) / 0.000 / 0.287 & 1 / 0.696 & 1 / 0.538 & 1 / 0.662 & 1 / 0.596 & 1 / 0.432 & 1 / 0.554 \\
A: 2×2, M1, 93 × 30 & 187 & 0.567 & 0.656 (0.015) / 0.002 / 0.160 & 1 / 0.694 & 1 / 0.477 & 1 / 0.964 & 1 / 0.633 & 1 / 0.414 & 1 / 0.932 \\
C: 2×2 pooled, truth near the boundary & 3 & 0.325 & 0.661 (0.015) / 0.076 / 0.031 & 1 / 0.087 & 0.999 (0.001) / 0.057 & 1 / 0.059 & 0.996 (0.002) / 0.041 & 0.994 (0.002) / 0.039 & 0.999 (0.001) / 0.041 \\
\textbf{E: 2×2 pooled, Slater fails} & 3 & 0.271 & \textbf{0.267 (0.014) / 0.503} / 0.038 & 1 / 0.249 & 0.998 (0.001) / 0.207 & 1 / 0.251 & 0.990 (0.003) / 0.179 & \textbf{0.944 (0.007)} / 0.143 & 0.995 (0.002) / 0.194 \\
B: 3×3 pooled, loyalty P→P & 12 & 0.835 & 1 / 0.000 / 0.533 & 1 / 0.675 & 1 / 0.659 & 1 / 0.679 & 1 / 0.592 & 1 / 0.557 & 1 / 0.581 \\
B: 3×3 pooled, defection P→O & 12 & 0.635 & 1 / 0.000 / 0.348 & 1 / 0.447 & 1 / 0.431 & 1 / 0.445 & 1 / 0.399 & 1 / 0.369 & 1 / 0.395 \\
B: 3×3 pooled, demobilisation & 12 & 0.533 & 1 / 0.000 / 0.283 & 1 / 0.381 & 1 / 0.363 & 1 / 0.378 & 1 / 0.333 & 1 / 0.304 & 1 / 0.329 \\
B: 3×3 pooled, mobilisation N→P & 12 & 0.884 & 1 / 0.000 / 0.528 & 1 / 0.683 & 1 / 0.646 & 1 / 0.661 & 1 / 0.619 & 1 / 0.559 & 1 / 0.595 \\
\end{longtable}

Design E puts the truth on the boundary so that Slater fails. Compositions span \([0.02,0.98]\) with both ends observed. The mean rate of group \(b\) is exactly 0 at the right vertex and that of group \(w\) is exactly 0 at the left vertex. The lower and upper constraints on \(k\) then meet at the truth, so the identified set of the model-implied rate is a single point, and no \(k\) satisfies every constraint strictly. Design C is the earlier boundary design: the rate is close to 0 at one vertex, so Slater holds, but only just. The CLR interval for \(B\) with the Popoviciu allowance covered in all 1000 replications of every row. Its widths are 0.530, 0.472, 0.086, 0.221 and 0.651/0.428/0.363/0.641 (rows in order); full output in \path{phase3_montecarlo_output.txt}.

\hypertarget{pE-many-small-districts-r2-m1-22-10-polls-per-district}{%
\subsubsection{8.2 Many small districts (R2): M1, 2×2, 10 polls per district}\label{pE-many-small-districts-r2-m1-22-10-polls-per-district}}

\begin{longtable}[]{@{}
  >{\raggedright\arraybackslash\hspace{0pt}}p{(\columnwidth - 20\tabcolsep) * \real{0.0909}}
  >{\raggedright\arraybackslash\hspace{0pt}}p{(\columnwidth - 20\tabcolsep) * \real{0.0909}}
  >{\raggedright\arraybackslash\hspace{0pt}}p{(\columnwidth - 20\tabcolsep) * \real{0.0909}}
  >{\raggedright\arraybackslash\hspace{0pt}}p{(\columnwidth - 20\tabcolsep) * \real{0.0909}}
  >{\raggedright\arraybackslash\hspace{0pt}}p{(\columnwidth - 20\tabcolsep) * \real{0.0909}}
  >{\raggedright\arraybackslash\hspace{0pt}}p{(\columnwidth - 20\tabcolsep) * \real{0.0909}}
  >{\raggedright\arraybackslash\hspace{0pt}}p{(\columnwidth - 20\tabcolsep) * \real{0.0909}}
  >{\raggedright\arraybackslash\hspace{0pt}}p{(\columnwidth - 20\tabcolsep) * \real{0.0909}}
  >{\raggedright\arraybackslash\hspace{0pt}}p{(\columnwidth - 20\tabcolsep) * \real{0.0909}}
  >{\raggedright\arraybackslash\hspace{0pt}}p{(\columnwidth - 20\tabcolsep) * \real{0.0909}}
  >{\raggedright\arraybackslash\hspace{0pt}}p{(\columnwidth - 20\tabcolsep) * \real{0.0909}}@{}}
\toprule\noalign{}
\begin{minipage}[b]{\linewidth}\raggedright\hspace{0pt}Districts
\end{minipage} & \begin{minipage}[b]{\linewidth}\raggedright\hspace{0pt}\(q\)
\end{minipage} & \begin{minipage}[b]{\linewidth}\raggedright\hspace{0pt}DD
\end{minipage} & \begin{minipage}[b]{\linewidth}\raggedright\hspace{0pt}Plug-in: cover / empty
\end{minipage} & \begin{minipage}[b]{\linewidth}\raggedright\hspace{0pt}CLR, \(\tilde B\)
\end{minipage} & \begin{minipage}[b]{\linewidth}\raggedright\hspace{0pt}Projection, \(\tilde B\)
\end{minipage} & \begin{minipage}[b]{\linewidth}\raggedright\hspace{0pt}Thm P3.8, \(\tilde B\)
\end{minipage} & \begin{minipage}[b]{\linewidth}\raggedright\hspace{0pt}CLR, \(B\)
\end{minipage} & \begin{minipage}[b]{\linewidth}\raggedright\hspace{0pt}Projection, \(B\)
\end{minipage} & \begin{minipage}[b]{\linewidth}\raggedright\hspace{0pt}Thm P3.8, \(B\)
\end{minipage} & \begin{minipage}[b]{\linewidth}\raggedright\hspace{0pt}Conservative, \(B\)
\end{minipage} \\
\midrule\noalign{}
\endhead
\bottomrule\noalign{}
\endlastfoot
30 & 61 & 0.571 & 0.305 (0.015) / 0.241 & 0.999 (0.001) / 0.557 & 1 / 0.967 & 1 / 1.000 & 1 / 0.746 & 1 / 1.000 & 1 / 1.000 & 1 / 1.000 \\
93 & 187 & 0.566 & 0.042 (0.006) / 0.463 & 0.995 (0.002) / 0.467 & 1 / 1.000 & 1 / 1.000 & 0.998 (0.001) / 0.575 & 1 / 1.000 & 1 / 1.000 & 1 / 0.930 \\
200 & 401 & 0.566 & 0.005 (0.002) / 0.717 & 0.996 (0.002) / 0.416 & 1 / 1.000 & 1 / 1.000 & 1 / 0.491 & 1 / 1.000 & 1 / 1.000 & 1 / 0.890 \\
\end{longtable}

Theorem P3.8 also covered in all 1000 replications of designs A 10 × 93 and 93 × 30, with median clipped width 1.000 in both, for \(\tilde B\) and \(B\).

\hypertarget{pE-riding-shocks-m0-against-m1-22-20-districts-50-polls}{%
\subsubsection{8.3 Riding shocks: M0 against M1, 2×2, 20 districts × 50 polls}\label{pE-riding-shocks-m0-against-m1-22-20-districts-50-polls}}

The district shifts are added to both groups' rates, and depend on the district's mean composition \(m_d\):

\begin{itemize}
\tightlist
\item
  \textbf{S0:} none;
\item
  \textbf{S1:} affine in \(m_d\), plus noise (as in design A);
\item
  \textbf{S2:} a common shock \(0.25\cos(2\pi m_d)\), non-affine in \(m_d\): the reviewer's common-shock example.
\end{itemize}

\begin{longtable}[]{@{}
  >{\raggedright\arraybackslash\hspace{0pt}}p{(\columnwidth - 16\tabcolsep) * \real{0.1111}}
  >{\raggedright\arraybackslash\hspace{0pt}}p{(\columnwidth - 16\tabcolsep) * \real{0.1111}}
  >{\raggedright\arraybackslash\hspace{0pt}}p{(\columnwidth - 16\tabcolsep) * \real{0.1111}}
  >{\raggedright\arraybackslash\hspace{0pt}}p{(\columnwidth - 16\tabcolsep) * \real{0.1111}}
  >{\raggedright\arraybackslash\hspace{0pt}}p{(\columnwidth - 16\tabcolsep) * \real{0.1111}}
  >{\raggedright\arraybackslash\hspace{0pt}}p{(\columnwidth - 16\tabcolsep) * \real{0.1111}}
  >{\raggedright\arraybackslash\hspace{0pt}}p{(\columnwidth - 16\tabcolsep) * \real{0.1111}}
  >{\raggedright\arraybackslash\hspace{0pt}}p{(\columnwidth - 16\tabcolsep) * \real{0.1111}}
  >{\raggedright\arraybackslash\hspace{0pt}}p{(\columnwidth - 16\tabcolsep) * \real{0.1111}}@{}}
\toprule\noalign{}
\begin{minipage}[b]{\linewidth}\raggedright\hspace{0pt}Design
\end{minipage} & \begin{minipage}[b]{\linewidth}\raggedright\hspace{0pt}Fit
\end{minipage} & \begin{minipage}[b]{\linewidth}\raggedright\hspace{0pt}Plug-in: cover / empty
\end{minipage} & \begin{minipage}[b]{\linewidth}\raggedright\hspace{0pt}CLR, \(\tilde B\): cover / empty / width
\end{minipage} & \begin{minipage}[b]{\linewidth}\raggedright\hspace{0pt}Projection, \(\tilde B\)
\end{minipage} & \begin{minipage}[b]{\linewidth}\raggedright\hspace{0pt}Conservative, \(\tilde B\)
\end{minipage} & \begin{minipage}[b]{\linewidth}\raggedright\hspace{0pt}CLR, \(B\)
\end{minipage} & \begin{minipage}[b]{\linewidth}\raggedright\hspace{0pt}Projection, \(B\)
\end{minipage} & \begin{minipage}[b]{\linewidth}\raggedright\hspace{0pt}Pooling test: rejection at 5\%
\end{minipage} \\
\midrule\noalign{}
\endhead
\bottomrule\noalign{}
\endlastfoot
S0 no shifts & M0 & 1 / 0.000 & 1 / 0.000 / 0.329 & 1 / 0.336 & 1 / 0.342 & 1 / 0.429 & 1 / 0.436 & \\
& M1 & 0.999 (0.001) / 0.000 & 1 / 0.000 / 0.468 & 1 / 0.704 & 1 / 0.667 & 1 / 0.571 & 1 / 0.808 & \textbf{0.090 (0.009)} (size; df 38) \\
S1 affine shifts & M0 & 1 / 0.000 & 1 / 0.000 / 0.335 & 1 / 0.344 & 1 / 0.351 & 1 / 0.435 & 1 / 0.444 & \\
& M1 & 0.938 (0.008) / 0.000 & 1 / 0.000 / 0.437 & 1 / 0.672 & 1 / 0.638 & 1 / 0.539 & 1 / 0.773 & 1.000 (0.000) \\
S2 common non-affine shock & M0 & \textbf{0.002 (0.001) / 0.987} & \textbf{0.009 (0.003) / 0.963} / 0.023 & \textbf{0.013 (0.004)}, empty 0.950 & \textbf{0.015 (0.004)}, empty 0.964 & \textbf{0.155 (0.011)}, empty 0.585 & \textbf{0.206 (0.013)}, empty 0.478 & \\
& M1 & 0.997 (0.002) / 0.002 & 1 / 0.000 / 0.328 & 1 / 0.509 & 1 / 0.541 & 1 / 0.429 & 1 / 0.611 & 1.000 (0.000) \\
\end{longtable}

\hypertarget{pE-reading-the-results}{%
\subsubsection{8.4 Reading the results}\label{pE-reading-the-results}}

\begin{enumerate}
\def\labelenumi{\arabic{enumi}.}
\tightlist
\item
  \textbf{No proved interval under-covers.}

  \begin{itemize}
  \tightlist
  \item
    Where its theorem applies, the projection interval covers at least 0.995 (design E, MC SE 0.002) and the conservative interval at least 0.990 (design E, SE 0.003).
  \item
    Theorem P3.8 covers in every replication of every design where it was computed, including R2.
  \item
    The only figures below 0.95 for a correctly specified model are the plug-in (not a CI) and CLR for \(\tilde B\) in design E.
  \end{itemize}
\item
  \textbf{CLR (a simulation finding) covers at least 0.994 in every design where Slater holds}, including R2 with 200 districts (0.995--0.999 for \(\tilde B\), 0.998--1 for \(B\)).

  \begin{itemize}
  \tightlist
  \item
    \textbf{At the Slater failure (design E), CLR for \(\tilde B\) covers 0.944 (SE 0.007).} This is within one standard error of the nominal 0.95, so it is not evidence of under-coverage.
  \item
    It is the informative case: the identified set is a point, so the interval has no slack, and coverage falls from near 1 to the nominal level. The 100\% elsewhere reflects truths in the interior of the identified set (§4).
  \item
    With the residual allowance (target \(B\)) CLR covers 0.998 in design E.
  \end{itemize}
\item
  \textbf{The plug-in fails where Theorem P3.4 predicts, and badly at the Slater failure.}

  \begin{itemize}
  \tightlist
  \item
    In R2 its coverage falls with the number of districts: 0.305, 0.042 and 0.005 at 30, 93 and 200 districts. It is empty 24\%, 46\% and 72\% of the time.
  \item
    In design E it covers 0.267 and is empty half the time: with a point identified set, any sampling noise empties the plug-in.
  \end{itemize}
\item
  \textbf{The projection interval is useful only when \(q\) is small.} Its critical value grows like \(\sqrt q\): 2.8 at \(q=3\), 4.6 at \(q=12\), 14.8 at \(q=187\) and 21.2 at \(q=401\), while CLR's is about 2--3.

  \begin{itemize}
  \tightlist
  \item
    Pooled designs (\(q=3\), 12): it is only 0.01--0.03 wider than CLR.
  \item
    M1 with 10 districts (\(q=21\)): 0.554 against CLR's 0.432 for \(\tilde B\).
  \item
    M1 with 93 districts: 0.93 at 30 polls per district, and \([0,1]\) at 10.
  \item
    So in grouped designs with many districts CLR is the only informative interval, and its support there is simulation only (§0.5).
  \end{itemize}
\item
  \textbf{Theorem P3.8 remains uninformative.} Its clipped width is 1.000 in every M1 design, including 10 districts × 93 polls, where CLR is 0.43 and DD is 0.57. This confirms §5.1 at 1000 replications.
\item
  \textbf{In R2, CLR narrows as districts are added.} For \(\tilde B\) its width is 0.557, 0.467 and 0.416 at 30, 93 and 200 districts; for \(B\), 0.746, 0.575 and 0.491. With 200 districts of 10 polls it is narrower than DD (0.566). The earlier statement that ``with 10 units per district the valid CI is no narrower than DD'' holds at 93 districts but not at 200.
\item
  \textbf{Riding shocks (R-37(i)).}

  \begin{itemize}
  \tightlist
  \item
    \textbf{Common non-affine shock (S2):} every M0 interval fails, for \(B\) and for \(\tilde B\). CLR for \(\tilde B\) covers 0.9\%, and is empty (ends crossing, so M0 is rejected) 96\% of the time. For \(B\), CLR covers 16\% and projection 21\% (the residual allowance widens them), and the conservative interval 4\%; again most failures are empty intervals. M1 covers in every replication, and the pooling test rejects in every replication.
  \item
    \textbf{Affine shifts (S1):} the pooling test rejects in every replication, because M0 is false for the observed design (P3.6(c)), but M0's intervals still cover. Shifts affine in composition change the pooled conditional mean only through terms the pooled model absorbs (2b Prop. D4(b)), approximately, so M0's bias is small here. A rejection of the pooling test therefore says that M0 is wrong, not that its intervals are.
  \item
    \textbf{Size of the pooling test (S0): 0.090 (SE 0.009) at nominal 0.05} (this Monte Carlo, 1000 replications, χ² calibration; the separate 2000-replication run in \path{benchmarks/pooling_test_size.R} gives 0.078 for χ², and 0.043 with the wild bootstrap now used by the package, P3.6 calibration table), with 38 restrictions and 50 polls per district. The HC3 Wald statistic has mean 39.7 and variance 100, against 38 and 76 for \(\chi^2_{38}\). The F form (\(W/\text{df}\) against \(F_{38,959}\)) gives 0.079. This is consistent with P3.6, whose \(\chi^2\) calibration is established for fixed dimension only, but the test over-rejects in a moderate-dimension R1 design. It errs towards M1, the safer rung. \emph{(Follow-up, 0.3.0: the restricted wild bootstrap rejects 0.043 (SE 0.005) here; P3.6, calibration paragraph.)}
  \end{itemize}
\item
  \textbf{The residual term is a large part of the width.} Targeting \(\tilde B\) instead of \(B\) narrows CLR by 0.06--0.10 in the 3×3 design and 0.06--0.11 in the 2×2 M1 designs with 30 or more polls per district; with 10 polls per district, by 0.08--0.19.
\item
  \textbf{Bhatia--Davis versus Popoviciu.} In the boundary design C, Bhatia--Davis narrows CLR for \(B\) from 0.086 to 0.057. In the Slater design E it narrows it from 0.221 to 0.207. Elsewhere the plausible range of rates includes 0.5, the allowance equals Popoviciu's, and the interval is up to 0.015 wider (the α share spent on the variance bound). It remains the package default.
\item
  \textbf{Real data} (NZ benchmark, \path{benchmarks/results/ei_datasets_nz_summary.md}; run with eiPartialID 0.2.0 during WS6, before the 0.3.x changes to the tests and the \texttt{proj\_actual} variance bound, which can move those columns slightly).

  \begin{itemize}
  \tightlist
  \item
    Per electorate, with about 64 polling places each: M0's CLR and projection intervals covered all 4302 cells, at about DD width.
  \item
    Pooling each election's 62--69 electorates and fitting M1 with electorates as groups: the M1 test (against M2) rejects M1 in all seven elections (\(p\approx0\)), and the pooling test rejects M0 against M1 (\(p\approx0\)). The plug-in M1 estimate is empty in all seven elections, but that is a sample diagnostic: the emptiness test does not reject M1 in any election (it rejects pooled M0 in 2002 and 2020 only).
  \item
    There, CLR intervals are rejected (ends crossing) in 21\% of cells and miss the truth in 8\%. The projection interval never misses, but it is as wide as DD.
  \item
    So where the model is wrong, the tests flag it, and a non-empty CLR interval is not reassurance.
  \end{itemize}
\end{enumerate}

\begin{center}\rule{0.5\linewidth}{0.5pt}\end{center}

\hypertarget{pE-practical-recipe-input-to-phase-4}{%
\subsection{9. Practical recipe (input to Phase 4)}\label{pE-practical-recipe-input-to-phase-4}}

\begin{enumerate}
\def\labelenumi{\arabic{enumi}.}
\tightlist
\item
  Fit the rung's regression (M0 / M1 / M2) on \(C-1\) columns with HC3 covariance (no clustering, §0.4), and the sample hulls.
\item
  Build the LP (2a Prop. 7 / 2b Thm 2) and report the plug-in bound, labelled as an estimate of the identified set, not a confidence interval.
\item
  Report the CLR confidence interval for \(\tilde B\) (default) and for \(B\), labelled algorithmic (Conjecture P3.5(b3); for \(B\) its Bhatia--Davis allowance is P3.5(c2)). Report the projection interval (P3.5(b2)(ii), with \(q=\dim\beta\) or an exact reduction) alongside it as the interval with an asymptotic guarantee; for \(B\) its Bhatia--Davis allowance must take \(\bar m\) from the projection interval on the rung alone (P3.5(c1)), or use Popoviciu. Intersect both with Duncan--Davis, which is valid because \(B\in DD\) always.
\item
  Theorem P3.8 gives a finite-sample guarantee in every regime. In the designs tried, its width is that of the logical range, so it is not yet worth reporting by default. With many small districts (R2), say that the CLR interval's coverage there rests on simulation (§0.5).
\item
  Report the emptiness test and, for M1, the test of its testable part (P3.6(b)) and the pooling test (P3.6(c)). The \(\chi^2\) version of the pooling test over-rejects with many restrictions (9\% at nominal 5\% with 38 restrictions, §8.4); since 0.3.0 the package calibrates both tests by a restricted wild bootstrap (4.3\% and 5.4\% in design S0; P3.6). A rejection says M0 is wrong for this design, not that its intervals fail (§8.3, S1).
\item
  Report the sensitivity curves in \(\eta\) (and \(\delta\)), with the breakdown value of the headline conclusion.
\item
  With few units per district, apply P3.7 selection, or prefer M1(\(\delta\)) or coarser groups, and say so.
\end{enumerate}

\hypertarget{pE-open-items}{%
\subsection{10. Open items}\label{pE-open-items}}

\begin{itemize}
\tightlist
\item
  \emph{Closed after review (Batch B, Nov 2026):} (1) P3.5(b2)(ii) now requires a positive definite \(\Sigma\) or a reduced full-rank parameterisation; a singular \(\hat V\) in the package's projection interval needs a guard (rank check, or \(q=\dim\beta\) with a warning). (2) The residual step of P3.5(a) is completed by the Chebyshev/Lindeberg subsequence argument. Both repairs leave every default and every allowance unchanged.
\item
  \emph{Added after review (round 2, Nov 2026):} uniform (in the design sequence) limit theory for the plug-in endpoint when candidates are within \(O(n^{-1/2})\) of binding without a convergent localisation (Proposition P3.4′ covers convergent localised slacks).
\item
  Conjecture P3.5(b3): validity of collecting dual vertices by simulation, with \(k_p\) computed over the collected set. A proved reduction to (b1) is given; recovery and the \(k_p\) effect are open.
\item
  Boundary cases where Slater fails (implied rates exactly 0 at a vertex): rate and limit of the endpoint estimator.
\item
  The rate of the hull error \(\varphi(\beta_0;\hat A)-\varphi(\beta_0;A)\), and inference for the population-hull endpoint.
\item
  Inference for the superpopulation target \(B^*\), which needs design-sampling variation.
\item
  Uniform validity near point identification for the CLR version (Stoye-type adaptation).
\item
  Asymptotic theory for M1/M2 in R3 (high-dimensional Gaussian approximation; §0.5) and in R2.
\item
  \textbf{Informative finite-sample inference in R2.} Theorem P3.8 is valid but uninformative in the designs tried. Candidates:

  \begin{itemize}
  \tightlist
  \item
    bound the endpoint directly rather than every constraint, e.g.~a finite-sample version of (b2) with a data-dependent dual;
  \item
    variance-adaptive bounds (empirical Bernstein);
  \item
    Bhatia--Davis;
  \item
    uneven \(\alpha\) splits.
  \end{itemize}
\item
  CLR when Slater fails: coverage fell to the nominal level (0.944, SE 0.007) in design E of §8, as expected without slack. Whether it holds uniformly near point identification is the Stoye-type question above.
\item
  Calibration of the pooling and M1 tests in moderate and high dimension. The HC3 Wald \(\chi^2\) test over-rejects (9\% at 5\% with 38 restrictions, §8.4). \emph{Partly addressed (0.3.0):} both tests now use a restricted wild bootstrap, which is close to nominal in design S0 (P3.6, calibration paragraph). Open: a proof for a growing number of restrictions.
\item
  Inference under R2 with a hierarchical model for the district shifts.
\item
  Dependence within districts beyond the district effect (spillovers; §0.4).
\item
  Joint relaxation for population turnover.
\item
  Coverage of the plug-in under misspecification, with the menu's dials.
\end{itemize}
